%% =============================================================================
%% SWR_v2_ger.tex
%% Synthetische Weltbild-Rekonstruktion: Ein LLM-gest\"utztes Verfahren
%% zur Analyse latenter \"Uberzeugungssysteme
%% Autor: Lukas Geiger, Unabh\"angiger Forscher, Bernau
%% Template: interdisciplinary (Methoden)
%% KI-Disclosure: L5 (Extensive)
%% =============================================================================
%%
%% Kompilieren mit: pdflatex -> bibtex -> pdflatex -> pdflatex
%% =============================================================================

\documentclass[11pt,a4paper]{article}

%% ---------------------------------------------------------------------------
%% Kodierung und Sprache
%% ---------------------------------------------------------------------------
\usepackage[utf8]{inputenc}
\usepackage[T1]{fontenc}
\usepackage[ngerman]{babel}

%% ---------------------------------------------------------------------------
%% Seitenlayout
%% ---------------------------------------------------------------------------
\usepackage[
  a4paper,
  left=2.5cm,
  right=2.5cm,
  top=2.5cm,
  bottom=2.5cm
]{geometry}

\usepackage{setspace}
\onehalfspacing

%% ---------------------------------------------------------------------------
%% Schriften
%% ---------------------------------------------------------------------------
\usepackage{mathptmx}
\usepackage[scaled=0.92]{helvet}

\usepackage{titlesec}
\titleformat{\section}
  {\normalfont\Large\bfseries\sffamily}{\thesection.}{0.5em}{}
\titleformat{\subsection}
  {\normalfont\large\bfseries\sffamily}{\thesubsection}{0.5em}{}
\titleformat{\subsubsection}
  {\normalfont\normalsize\bfseries\sffamily}{\thesubsubsection}{0.5em}{}

%% ---------------------------------------------------------------------------
%% Zitierweise
%% ---------------------------------------------------------------------------
\usepackage[round,authoryear,semicolon]{natbib}
\bibliographystyle{plainnat}

%% ---------------------------------------------------------------------------
%% Pakete
%% ---------------------------------------------------------------------------
\usepackage{graphicx}
\usepackage{booktabs}
\usepackage{array}
\usepackage{tabularx}
\usepackage{amsmath,amssymb}
\usepackage{amsthm}
\usepackage{enumitem}
\usepackage{xcolor}
\usepackage{url}
\usepackage{longtable}
\usepackage[breaklinks=true]{hyperref}
\hypersetup{
  colorlinks  = true,
  linkcolor   = black,
  citecolor   = blue!60!black,
  urlcolor    = blue!70!black,
  pdfauthor   = {Lukas Geiger},
  pdftitle    = {Synthetische Weltbild-Rekonstruktion},
  pdfsubject  = {Methoden der qualitativen Sozialforschung},
}

%% ---------------------------------------------------------------------------
%% Umgebungen
%% ---------------------------------------------------------------------------
\theoremstyle{definition}
\newtheorem{definition}{Definition}[section]

\theoremstyle{plain}
\newtheorem{these}{These}[section]

\theoremstyle{remark}
\newtheorem{anmerkung}{Anmerkung}[section]

%% ---------------------------------------------------------------------------
%% Titelblock
%% ---------------------------------------------------------------------------
\title{%
  \sffamily\bfseries
  Synthetische Weltbild-Rekonstruktion:\\[0.3em]
  \large Ein LLM-gest\"utztes Verfahren zur Analyse\\
  latenter \"Uberzeugungssysteme%
}

\author{%
  Lukas Geiger\thanks{Korrespondenz: Lukas Geiger, Bernau, Deutschland. ORCID: \url{https://orcid.org/0009-0005-7296-1534}}\\[0.3em]
  \small Unabh\"angiger Forscher, Bernau
}

\date{M\"arz 2026}

%% ===========================================================================
\begin{document}
%% ===========================================================================

\maketitle
\thispagestyle{empty}

%% ---------------------------------------------------------------------------
%% Zusammenfassung
%% ---------------------------------------------------------------------------
\begin{abstract}
\noindent
Die vorliegende Arbeit stellt die \emph{Synthetische Weltbild-Rekonstruktion} (SWR) als neue sozialwissenschaftliche Methode vor. Die SWR formalisiert die Rekonstruktion latenter \"Uberzeugungssysteme als \emph{inverses Problem}: Gegeben sind die beobachtbaren Manifestationen eines Akteurs -- \"offentliche Aussagen und dokumentierte Handlungen --, gesucht ist das innere \"Uberzeugungssystem (Weltbild), das diese Manifestationen am koh\"arentesten erzeugen w\"urde. Ein gro\ss{}es Sprachmodell (Large Language Model, LLM) fungiert als Inversionsoperator und konstruiert eine \emph{fiktive synthetische Person}, deren Weltbild den bestm\"oglichen Fit f\"ur die Gesamtheit der Daten bietet. Das Verfahren operationalisiert diesen Grundgedanken durch ein standardisiertes F\"unfschritt-Prompt-Protokoll (Datenkorpus zusammenstellen, Synthese, Interview, Reflexion, Destillat), eine systematische Datengruppierung in Syntheseeinheiten entlang von sechs unabh\"angigen Variablen, sechs komplement\"are Operationalisierungen und einen mehrschichtigen Validierungsrahmen. Die theoretische Verankerung verbindet vier Traditionen: Webers Idealtyp, Booths \emph{implied author}, Inverse Reinforcement Learning und projektive Testverfahren. Die Methode verbindet qualitative Tiefe mit quantitativer Skalierbarkeit und zielt darauf ab, die systematische Rekonstruktion und den Vergleich latenter Weltbilder f\"ur gro\ss{}e Personenzahlen zu erm\"oglichen.

Die empirische Demonstration der Methode erfolgt in einem Companion Paper (DOI: 10.5281/zenodo.18736737), das die vollst\"andige SWR auf $N=100$ Personen der KI-Elite anwendet.

\medskip
\noindent\textbf{Schlagw\"orter:} Weltbild-Rekonstruktion, Large Language Models, qualitative Methoden, inverses Problem, Computational Social Science, fiktive synthetische Person
\end{abstract}

\newpage
\setcounter{page}{1}

%% ===========================================================================
%% 1. EINLEITUNG
%% ===========================================================================

\section{Einleitung}
\label{sec:einleitung}

\subsection{Problemstellung}
\label{sec:problemstellung}

Einflussreiche Akteure -- politische Entscheidungstr\"ager, Unternehmenslenkende, \"offentliche Intellektuelle -- handeln nicht in einem Vakuum rationaler Kalkulation.\footnote{Diese Arbeit verwendet durchgehend das redaktionelle \glqq wir\grqq{} gem\"a\ss{} akademischer Konvention. Der alleinige menschliche Autor ist Lukas Geiger; die KI-Beteiligung ist in der KI-Offenlegungserkl\"arung dokumentiert.} Ihren Entscheidungen, Aussagen und Strategien liegen implizite \"Uberzeugungssysteme zugrunde: grundlegende Annahmen \"uber die Natur des Menschen, die Struktur der Gesellschaft, die Rolle von Technologie, das Verh\"altnis von Individuum und Gemeinschaft. Diese \"Uberzeugungssysteme -- im Folgenden als \emph{Weltbilder} bezeichnet -- sind selten explizit artikuliert, wirken aber als kognitive Tiefenstrukturen, die sowohl die Problemwahrnehmung als auch den L\"osungsraum eines Akteurs determinieren \citep{Mannheim1929, BergerLuckmann1966}.

Die systematische Rekonstruktion solcher Weltbilder stellt die Sozialwissenschaften vor ein fundamentales methodisches Problem. Weltbilder sind \emph{latente Konstrukte}: Sie sind nicht direkt beobachtbar, sondern manifestieren sich nur indirekt -- in \"offentlichen Reden, Interviews, Policy-Dokumenten, unternehmerischen Entscheidungen, ver\"offentlichten Texten und beobachtbarem Handeln. Der Forscher steht vor der Aufgabe, aus einer heterogenen Menge solcher Manifestationen das zugrunde liegende \"Uberzeugungssystem zu inferieren -- ein Problem, das in seiner epistemologischen Struktur einem \emph{inversen Problem} entspricht: Gegeben sind die Outputs (beobachtbare Aussagen und Handlungen), gesucht ist der generierende Mechanismus (das Weltbild), der diese Outputs hervorbringt.

\subsection{Grenzen bestehender Verfahren}
\label{sec:grenzen}

Die Sozialwissenschaften verf\"ugen \"uber ein breites Repertoire an Methoden zur Analyse von Texten, Diskursen und Sinnstrukturen. F\"ur die hier formulierte Aufgabe -- die systematische Rekonstruktion latenter Weltbilder aus \"offentlich verf\"ugbaren Quellen f\"ur eine gro\ss{}e Zahl von Akteuren -- sto\ss{}en diese Verfahren jedoch an charakteristische Grenzen.

Die \emph{qualitative Inhaltsanalyse} nach \citet{Mayring2022} bietet ein strukturiertes, regelgeleitetes Vorgehen mit hoher intersubjektiver Nachvollziehbarkeit, skaliert aber schlecht und bietet keinen genuinen Mechanismus zur Synthese eines koh\"arenten Gesamtbilds. Die \emph{Kritische Diskursanalyse} \citep{Fairclough1995} erlaubt die Aufdeckung impliziter Annahmen und Machtstrukturen, ist aber in hohem Ma\ss{}e interpretationsabh\"angig \citep{Breeze2011} und bietet kein standardisiertes Aggregationsformat. Die \emph{Frame-Analyse} \citep{Goffman1974, Entman1993} identifiziert kognitive Deutungsrahmen, reduziert aber das \"Uberzeugungssystem auf einzelne, isolierte Frames. Die \emph{Narrativanalyse} \citep{Riessman2008} erfasst integrative Sinnstrukturen, ist aber einzelfallzentriert und bietet keine Aggregationsverfahren. \emph{Automatisierte Textanalyseverfahren} -- Topic Modeling \citep{BleiNgJordan2003}, Sentiment-Analyse, Word Embeddings -- bieten Skalierbarkeit, operieren aber an der sprachlichen Oberfl\"ache.

Zusammenfassend l\"asst sich feststellen: Kein bestehendes Verfahren kann gleichzeitig (a)~latente \"Uberzeugungssysteme als koh\"arente Ganzheiten rekonstruieren, (b)~heterogene Datenquellen -- sowohl sprachliche \"Au\ss{}erungen als auch beobachtbares Handeln -- integrieren, (c)~auf eine gro\ss{}e Zahl von Akteuren (100+) skalieren und (d)~strukturierte Gruppenaggregate bilden, die den systematischen Vergleich erm\"oglichen.

\subsection{LLMs als neue M\"oglichkeit}
\label{sec:llms}

Die j\"ungste Generation gro\ss{}er Sprachmodelle (Large Language Models, LLMs) er\"offnet eine grundlegend neue Perspektive auf dieses methodische Problem. LLMs verf\"ugen \"uber zwei Eigenschaften, die sie als qualitative Forschungsinstrumente potenziell geeignet machen: Erstens haben sie im Zuge ihres Trainings auf umfangreichen Textkorpora ein breites Weltwissen absorbiert. Zweitens sind sie in der Lage, aus fragmentarischen Informationen koh\"arente Synthesen zu generieren -- eine F\"ahigkeit, die der hermeneutischen Kompetenz menschlicher Forscher strukturell \"ahnelt, aber in einem Bruchteil der Zeit operiert. Diese Perspektive kn\"upft an eine wachsende Literatur an, die LLMs als Analysewerkzeuge f\"ur die Sozialwissenschaften konzeptualisiert \citep{Ziems2024, Bail2024, Argyle2023}. J\"ungere Studien belegen die Eignung von LLMs f\"ur qualitative Kodieraufgaben: \citet{Tai2024} fanden, dass LLMs deduktive Kodierung mit einer Leistung vergleichbar zu menschlichen Kodierern durchf\"uhren k\"onnen; \citet{Barros2024} erstellten ein systematisches Mapping LLM-basierter qualitativer Forschung und dokumentierten einen starken Anstieg der Publikationst\"atigkeit 2023--2024; und Pilotstudien zur LLM-gest\"utzten Thematischen Analyse zeigten, dass Evaluatoren LLM-generierte Codes in einer Mehrzahl der F\"alle bevorzugten \citep{Wang2025}. Die SWR unterscheidet sich von diesen Ans\"atzen in einem entscheidenden Punkt: Statt bestehende Kategorien zu kodieren, f\"uhrt sie eine \emph{synthetische Integration} durch -- die Konstruktion koh\"arenter Weltbild-Modelle aus fragmentarischen Daten.

Zugleich bringen LLMs spezifische methodische Herausforderungen mit sich: das \emph{Kontaminationsrisiko} (die Modelle wurden auf denselben Quellen trainiert, die analysiert werden), den \emph{Koh\"arenzbias} (die Tendenz, widerspr\"uchliche Daten zu gl\"atten) und die \emph{Stochastizit\"at} der Modellausgaben.

\subsection{Die SWR-Methode: \"Uberblick}
\label{sec:ueberblick}

Die vorliegende Arbeit stellt eine Methode vor, die diese L\"ucke zu schlie\ss{}en sucht: die \emph{Synthetische Weltbild-Rekonstruktion} (SWR). Die SWR formuliert die Rekonstruktion impliziter Weltbilder als inverses Problem und nutzt ein LLM als Inversionsoperator. Das Verfahren erzeugt f\"ur systematisch definierte Datenb\"undel (\emph{Syntheseeinheiten}) jeweils eine \emph{fiktive synthetische Person}, deren Weltbild die Gesamtheit der pr\"asentierten Daten am koh\"arentesten erkl\"art. Sechs komplement\"are Operationalisierungen erschlie\ss{}en die latenten Konstrukte aus verschiedenen methodischen Perspektiven; ein mehrschichtiger Validierungsrahmen kontrolliert die spezifischen Risiken LLM-gest\"utzter Forschung.

\subsection{Beitrag}
\label{sec:beitrag}

Die vorliegende Arbeit leistet drei zentrale Beitr\"age:

\emph{Erstens} stellt sie mit der SWR, nach unserem Kenntnisstand, eine der ersten systematischen Methoden zur LLM-gest\"utzten Rekonstruktion latenter \"Uberzeugungssysteme vor. Die Methode formalisiert den Rekonstruktionsprozess als inverses Problem, operationalisiert Weltbilder \"uber sechs komplement\"are Zug\"ange und standardisiert die Generierung durch strukturierte Prompt-Sequenzen.

\emph{Zweitens} entwickelt die Arbeit einen sechsschichtigen Validierungsrahmen (V1--V6) f\"ur LLM-gest\"utzte qualitative Forschung, der Inter-LLM-Reliabilit\"at, Prompt-Sensitivit\"atsanalyse, dimensionale Quantifizierung, Identifikationstest, Run-Konvergenz-Messung und Adversarial Probing umfasst.

\emph{Drittens} formalisiert sie eine vierstufige Variablenarchitektur (abh\"angige Variablen $\rightarrow$ Operationalisierungen $\rightarrow$ Messvariablen $\rightarrow$ unabh\"angige Variablen) und ein System von neun Reduktionskriterien, die den kombinatorischen Datenraum auf eine handhabbare Zahl von Syntheseeinheiten reduzieren.

Der Beitrag dieser Arbeit liegt nicht in der Nutzung von LLMs zur Textanalyse -- dies ist bekannt --, sondern in der systematischen Methodisierung dieser Nutzung: ein reproduzierbares F\"unfschritt-Protokoll, ein sechsschichtiger Validierungsrahmen und eine formale Syntheseeinheiten-Systematik, die zusammen eine rigorose qualitative Methode konstituieren.


%% ===========================================================================
%% 2. THEORETISCHER RAHMEN
%% ===========================================================================

\section{Theoretischer Rahmen}
\label{sec:theorie}

\subsection{Weltbild als latentes Konstrukt}
\label{sec:weltbild}

Der Begriff \emph{Weltbild} (engl.\ \emph{worldview}, \emph{Weltanschauung}) bezeichnet in der hier vorgeschlagenen Verwendung die Gesamtheit der grundlegenden \"Uberzeugungen eines Akteurs \"uber das Selbst, die Welt und die Menschheit. Ein Weltbild umfasst ontologische Annahmen, epistemische \"Uberzeugungen, axiologische Orientierungen sowie praxeologische Dispositionen. Der Begriff ist bewusst umfassender gew\"ahlt als verwandte Termini wie \emph{Einstellung}, \emph{Werthaltung} oder \emph{Ideologie}, da er das \"Uberzeugungssystem als integrale Ganzheit zu erfassen beansprucht.

Diese Begriffsbestimmung steht in einer langen ideengeschichtlichen Tradition. \citet{Dilthey1911} f\"uhrte den Weltanschauungsbegriff als Grundkategorie der Geisteswissenschaften ein. \citet{Jaspers1919} systematisierte Weltanschauungen als umfassende Orientierungsstrukturen. In der Wissenssoziologie verwandte \citet{Mannheim1929} den Begriff der Weltanschauung, um standortgebundene Denkstile zu charakterisieren.

Eine zentrale Eigenschaft von Weltbildern ist ihre \emph{Latenz}: Sie sind nicht direkt beobachtbar. Kein Akteur verf\"ugt in der Regel \"uber eine vollst\"andige, explizite Artikulation seines Weltbilds. Stattdessen manifestieren sich Weltbilder indirekt -- in sprachlichen \"Au\ss{}erungen, in Handlungsentscheidungen, in dem, was als selbstverst\"andlich vorausgesetzt und nicht eigens thematisiert wird. Die Latenz hat eine doppelte Ursache: Viele grundlegende \"Uberzeugungen sind \emph{pr\"a-reflexiv} -- sie operieren als \emph{Doxa} im Sinne \citeauthor{Bourdieu1977}s~(\citeyear{Bourdieu1977}) --, und Weltbilder sind \emph{strategisch verdeckt}: Akteure artikulieren ihre \"Uberzeugungen selektiv \citep{Goffman1959}. Die methodische Konsequenz: Die Rekonstruktion erfordert einen \emph{abduktiven Inferenzschritt} im Sinne von \citet{Peirce1903} -- einen Schluss auf die beste Erkl\"arung.

Tabelle~\ref{tab:konstrukte} grenzt den Weltbild-Begriff von verwandten Konstrukten ab.

\begin{table}[htbp]
  \centering
  \caption{Abgrenzung des Weltbild-Begriffs von verwandten Konstrukten.}
  \label{tab:konstrukte}
  \begin{tabularx}{\textwidth}{l X l l l}
    \toprule
    \textbf{Konstrukt} & \textbf{Bezugseinheit} & \textbf{Umfang} & \textbf{Funktion} & \textbf{Latenz} \\
    \midrule
    Weltbild (SWR) & Individuum & Umfassend & Deskriptiv & Latent \\
    Ideologie & Soziale Gruppe & Umfassend & Herrschaftsbezogen & Latent \\
    Habitus & Individuum/Klasse & Inkl.\ K\"orper & Reproduktionsbez. & Latent \\
    Frame & Thema/Situation & Partikul\"ar & Aufmerksamkeit & Manifest/Lat. \\
    Mental Model & Individuum & Dom\"anenspez. & Kognitiv-funkt. & Latent \\
    \bottomrule
  \end{tabularx}
\end{table}


\subsection{\"Uberzeugungssystem-Forschung: Historischer und methodischer Kontext}
\label{sec:ansaetze}

Die Erforschung von \"Uberzeugungssystemen verf\"ugt \"uber eine reiche methodische Tradition, auf der die SWR aufbaut. \citet{Converse1964} zeigte in seiner wegweisenden Analyse der Massen\"offentlichkeit, dass politische \"Uberzeugungssysteme hochvariable Grade an \emph{Constraint} aufweisen -- die logische Verkn\"upfung zwischen einzelnen \"Uberzeugungen variiert dramatisch \"uber Bildungsschichten. Dieser Befund unterstreicht die Herausforderung, vor der die SWR steht: Der \glqq Koh\"arenzbias\grqq{} von LLMs (Sektion~\ref{sec:diskussion}) k\"onnte st\"arker verkn\"upfte Rekonstruktionen erzeugen als die empirische Realit\"at rechtfertigt, insbesondere bei Nicht-Eliten-Subjekten. \citeauthor{Axelrod1976}s~(\citeyear{Axelrod1976}) Technik des Cognitive Mapping repr\"asentiert einen fr\"uhen systematischen Versuch, kausale \"Uberzeugungsstrukturen aus verbalen Daten zu rekonstruieren; die SWR l\"asst sich als computationelle Erweiterung dieser Tradition verstehen. Die Q-Methodologie \citep{Stephenson1953,Watts2012} bietet einen besonders relevanten Vergleich: Wie die SWR zielt sie darauf ab, koh\"arente Muster von Subjektivit\"at zu identifizieren, jedoch durch Faktorenanalyse von Forced-Rank-Sortierungen statt durch LLM-vermittelte Synthese. Die SWR unterscheidet sich von der Q-Methodologie in drei Hinsichten: (i)~sie operiert auf naturalistischen Daten statt auf forscher-designten Q-Sorts; (ii)~sie produziert narrative Rekonstruktionen statt numerischer Faktorladungen; und (iii)~sie ist auf gro\ss{}e Personenzahlen skalierbar, ohne deren aktive Teilnahme zu erfordern. Ein systematischer Vergleich von SWR und Q-Methodologie am selben Datensatz w\"are eine wertvolle Validierungsübung.


\subsection{LLMs als qualitative Forschungsinstrumente}
\label{sec:llm-instrumente}

Die f\"ur die SWR entscheidende Eigenschaft von LLMs l\"asst sich als \emph{synthetische Integrationsf\"ahigkeit} bezeichnen: die F\"ahigkeit, aus fragmentarischen, heterogenen und verteilt vorliegenden Informationen koh\"arente Ganzheiten zu erzeugen. Diese F\"ahigkeit weist eine bemerkenswerte strukturelle Analogie zur \emph{hermeneutischen Kompetenz} auf: Der hermeneutische Zirkel -- das wechselseitige Erhellen von Teil und Ganzem \citep{Gadamer1960} -- findet in der Arbeitsweise von LLMs ein operationales Analogon.

Mehrere Arbeiten bilden den unmittelbaren Referenzrahmen: \citet{Ziems2024} zeigen, dass LLMs implizite Bedeutungen erfassen k\"onnen, die sich statistischen Oberfl\"achenanalysen entziehen. \citet{Bail2024} argumentiert, dass LLMs komplexe sozialwissenschaftliche Konstrukte messen k\"onnten. \citet{Argyle2023} demonstrieren, dass LLMs die Perspektiven heterogener sozialer Gruppen simulieren k\"onnen. \citet{Park2023} zeigen, dass generative Agenten plausibles soziales Verhalten erzeugen.

Diese strukturelle Analogie l\"asst sich durch die Perspektive der \emph{Computational Hermeneutics} sch\"arfen. \citet{Ramberg2026} schlagen vor, dass generative KI-Systeme als \glqq Kontextmaschinen\grqq{} fungieren, die drei interpretative Herausforderungen adressieren m\"ussen, die jeder hermeneutischen T\"atigkeit inh\"arent sind: \emph{Situiertheit} (Bedeutung entsteht nur im Kontext), \emph{Pluralit\"at} (mehrere g\"ultige Interpretationen koexistieren) und \emph{Ambiguit\"at} (Interpretationen konfligieren nat\"urlicherweise). Jede dieser Herausforderungen bildet sich direkt auf das methodische Design der SWR ab: Situiertheit wird durch die Syntheseeinheiten-Architektur adressiert, Pluralit\"at durch mehrere unabh\"angige Runs und Inter-LLM-Vergleich, Ambiguit\"at durch den Meta-Reflexionsschritt. Fundamentaler noch bietet Gadamers Konzept der \emph{Horizontverschmelzung} \citep{Gadamer1960} -- die produktive Verschmelzung des Vorverst\"andnisses des Interpreten mit dem interpretierten Artefakt -- einen philosophischen Rahmen f\"ur das, was die SWR operationalisiert: Das LLM bringt seinen trainingsbasierten \glqq Horizont\grqq{} in Kontakt mit dem Datenkorpus und erzeugt ein synthetisches Weltbild, das weder reine Datenextraktion noch reine Modellprojektion ist, sondern eine hermeneutische Fusion beider. Diese Analogie hat jedoch Grenzen: LLMs fehlt die genuine \emph{Vorurteilsstruktur}, die Gadamer als konstitutiv f\"ur Verstehen betrachtet, und ihr \glqq Horizont\grqq{} ist ein statistisches Artefakt, keine gelebte Perspektive \citep{Hornby2024}.

Die synthetische Integrationsf\"ahigkeit ist zugleich die Quelle spezifischer methodischer Risiken: (1)~\emph{Kontaminationsrisiko} -- das LLM verf\"ugt \"uber Vorwissen \"uber die analysierte Person; (2)~\emph{Koh\"arenzbias} -- LLMs produzieren systematisch koh\"arentere Narrative als die Daten rechtfertigen; (3)~\emph{Stochastizit\"at} -- identische Inputs k\"onnen verschiedene Outputs erzeugen; (4)~\emph{Implizite Normativit\"at} -- Trainingsdaten und RLHF pr\"agen die Outputs. Die SWR adressiert jedes dieser Risiken durch dedizierte Kontrollmechanismen (Sektion~\ref{sec:qualitaet}).


%% ===========================================================================
%% 3. DIE METHODE
%% ===========================================================================

\section{Die Methode: Synthetische Weltbild-Rekonstruktion}
\label{sec:methode}

\subsection{Grundidee: Das Weltbild als inverses Problem}
\label{sec:inversion}

Die SWR formalisiert die Rekonstruktion latenter Weltbilder als \emph{inverses Problem} \citep{Tarantola2005}. Das Vorw\"artsproblem lautet: Wenn eine Person mit Weltbild~$W$ in Situation~$S$ agiert, welche Aussagen~$A$ und Handlungen~$H$ sind zu erwarten?
%
\begin{equation}
  W + S \;\longrightarrow\; \{A, H\}
  \label{eq:forward}
\end{equation}
%
Das inverse Problem kehrt die Inferenzrichtung um:
%
\begin{equation}
  \{A_1, \ldots, A_n\} + \{H_1, \ldots, H_m\} \;\longrightarrow\; W^*
  \label{eq:inverse}
\end{equation}
%
Dieses Problem ist \emph{schlecht gestellt} (ill-posed) im Sinne Hadamards \citep{Hadamard1923}: Die drei klassischen Wohlgestelltheitskriterien -- Existenz, Eindeutigkeit und stetige Abh\"angigkeit von den Daten -- sind jeweils verletzt. \emph{Existenz} ist nicht garantiert, weil dem beobachteten Verhalten m\"oglicherweise kein koh\"arentes Weltbild zugrunde liegt. \emph{Eindeutigkeit} scheitert, weil mehrere intern konsistente Weltbilder dieselbe Menge an Manifestationen erkl\"aren k\"onnen. \emph{Stabilit\"at} ist gef\"ahrdet, weil kleine \"Anderungen im Datenkorpus (Hinzuf\"ugen oder Entfernen einer einzelnen Rede) das rekonstruierte Weltbild substantiell verschieben k\"onnen. Die SWR beansprucht daher nicht, \emph{das} Weltbild einer Person zu identifizieren, sondern \emph{ein} Weltbild zu rekonstruieren, das den bestm\"oglichen Fit f\"ur die Gesamtheit der beobachteten Daten bietet -- eine Bestapproximation unter Unsicherheit. Die methodische Antwort spiegelt -- auf struktureller, nicht mathematischer Ebene -- bew\"ahrte Strategien der inversen Problemtheorie wider: Regularisierung durch die Koh\"arenzinstruktion (strukturell analog zur Tikhonov-Regularisierung, insofern sie inkoh\"arente L\"osungen bestraft, jedoch als nat\"urlichsprachliche Instruktion implementiert statt als mathematischer Strafterm), Varianzkontrolle durch mehrere unabh\"angige Runs (analog zu Ensemble-Methoden) und Triangulation durch kreuzmodale Vorhersage (analog zur Multi-Source-Inversion). Diese Analogien verdeutlichen die epistemologische Struktur der Designentscheidungen der SWR; sie implizieren nicht, dass die SWR die formalen Garantien mathematischer Inversionsverfahren erreicht.

\begin{anmerkung}[Zur mathematischen Notation]
Die Gleichungen~\eqref{eq:forward} und~\eqref{eq:inverse} dienen als \emph{konzeptionelle Notation} -- sie formalisieren die logische Struktur des inversen Problems, sind aber nicht im mathematischen Sinne operational berechenbar. Die \glqq Inversion\grqq{} wird vom LLM durch sein Prompt-Protokoll durchgef\"uhrt, nicht durch einen mathematischen Algorithmus. Die Notation dient der Pr\"azisierung der epistemologischen Struktur der Methode, in der Tradition der inversen Problemtheorie \citep{Tarantola2005}, in der die Problemformulierung der (oft approximativen) L\"osungsmethode vorausgeht.
\end{anmerkung}


\subsection{Theoretische Verankerung}
\label{sec:verankerung}

Die Struktur des inversen Problems verbindet die SWR mit vier bedeutenden intellektuellen Traditionen:

\emph{Webers Idealtyp.} Max \citeauthor{Weber1904}s~(\citeyear{Weber1904}) Idealtyp ist ein gedankliches Konstrukt, das durch einseitige Steigerung bestimmter Gesichtspunkte gewonnen wird. Die synthetische Person ist ein Idealtypus des Weltbilds: ein konstruiertes Modell, das die Koh\"arenz der beobachteten Daten maximiert, ohne das tats\"achliche innere Erleben abzubilden.

Zugleich bestehen wesentliche Disanalogien: W\"ahrend Weber den Idealtyp als bewusste, forscherkontrollierte Abstraktion konzipierte, wird die Synthetische Person durch ein probabilistisches Modell erzeugt, dessen interner Inferenzprozess nicht transparent ist. Der Forscher kontrolliert den Prompt, nicht die Synthese. Dies impliziert einen Trade-off zwischen Kontrolle und Skalierbarkeit, der in der epistemologischen Reflexion ber\"ucksichtigt werden muss.

\emph{Booths Implied Author.} \citeauthor{Booth1961}s~(\citeyear{Booth1961}) \emph{implied author} ist das aus der Gesamtheit eines Werkes rekonstruierte koh\"arente \"Uberzeugungssubjekt. So wie der implied author aus den Entscheidungen eines Textes r\"uckerschlossen wird, wird das SWR-Weltbild aus den Entscheidungen eines Akteurs r\"uckerschlossen.

\emph{Inverse Reinforcement Learning.} IRL \citep{NgRussell2000} fragt: Welche Zielfunktion m\"usste ein rationaler Agent maximieren, damit das beobachtete Verhalten als optimale Strategie erkl\"arbar wird? An die Stelle der Belohnungsfunktion tritt das Weltbild, an die Stelle des beobachteten Verhaltens die dokumentierten Aussagen und Handlungen.

\emph{Projektive Tests.} In der klinischen Psychologie werden projektive Verfahren (Rorschach, TAT) eingesetzt, um aus den \"Au\ss{}erungen und Reaktionen einer Person auf deren innere Dispositionen zu schlie\ss{}en. Die SWR teilt die Grundlogik des R\"uckschlusses von \"Au\ss{}erung auf Inneres, transponiert auf die Ebene \"offentlicher Kommunikation.

Tabelle~\ref{tab:analogien} fasst die disziplin\"aren Parallelen zusammen.

\begin{table}[htbp]
  \centering
  \caption{Disziplin\"are Parallelen zum inversen Ansatz der SWR.}
  \label{tab:analogien}
  \begin{tabularx}{\textwidth}{l l X}
    \toprule
    \textbf{Disziplin} & \textbf{Konzept} & \textbf{Parallele} \\
    \midrule
    Physik & Inverses Problem & Aus Messdaten das Modell rekonstruieren \\
    Psychologie & Projektive Tests & Aus \"Au\ss{}erungen auf Inneres schlie\ss{}en \\
    Literaturwiss. & Implied Author & Der Autor, den der Text impliziert \\
    Soziologie & Idealtyp (Weber) & Verdichtetes Konstrukt, kein Abbild \\
    Statistik & Latente Variablen & Unsichtbare Faktoren hinter Daten \\
    KI-Forschung & Inverse Reinforcement Learning & Aus Verhalten auf Belohnungsfunktion schlie\ss{}en \\
    \bottomrule
  \end{tabularx}
\end{table}


\subsection{Die fiktive synthetische Person}
\label{sec:person}

Das zentrale methodische Konstrukt der SWR ist die \emph{fiktive synthetische Person}: ein vom LLM erzeugtes, fiktives Individuum, dessen \"Uberzeugungssystem so konstruiert ist, dass es die Gesamtheit der pr\"asentierten Daten am koh\"arentesten erkl\"art. Die synthetische Person ist weder eine Simulation der realen Person noch ein Portr\"at; sie ist das Destillat -- das Weltbild in Reinform, befreit von biographischem Rauschen, PR-Strategien und situativen Verzerrungen. Sie ist die Antwort auf die Frage:

\begin{quote}
\emph{Welcher Mensch m\"usste existieren, damit all diese Aussagen und Handlungen aus einem koh\"arenten Weltbild hervorgehen?}
\end{quote}

Die Einf\"uhrung einer fiktiven Person erf\"ullt f\"unf zentrale methodische Funktionen:

\begin{enumerate}[leftmargin=*]
  \item \textbf{Kontaminationsvermeidung.} Wenn das LLM eine fiktive Person konstruiert, wird das Risiko reduziert, dass Vorwissen aus dem Trainingskorpus einflie\ss{}t.
  \item \textbf{Trennung von Weltbild und Biographie.} Die fiktive Person trennt das Denken von der biographischen Kontingenz -- ein Idealtypus des Denkens, kein Abbild des Lebens.
  \item \textbf{Aggregationsf\"ahigkeit.} Bei Gruppen-Syntheseeinheiten gibt es keine einzelne reale Person. Aber es gibt ein synthetisches Kollektiv-Weltbild, das sich als fiktive Person ausdr\"ucken l\"asst.
  \item \textbf{Standardisierte Vergleichbarkeit.} Alle synthetischen Personen werden mit identischem Prompt und Verfahren erzeugt -- wie Messwerte aus demselben Instrument.
  \item \textbf{Sichtbarmachung von Widerspr\"uchen.} Eine fiktive Person \emph{muss} die Widerspr\"uche in den Daten integrieren. Die Reflexion macht diese Spannungen explizit.
\end{enumerate}

Eine fundamentale epistemologische Frage bleibt bestehen: Rekonstruiert das LLM tats\"achlich ein latentes \"Uberzeugungssystem, oder produziert es lediglich eine plausibel klingende Textcollage, die der Leser nachtr\"aglich als koh\"arentes Weltbild interpretiert? Die SWR kann diese Frage nicht definitiv beantworten -- sie ist eine Instanz des breiteren Problems, ob LLM-Outputs Verstehen konstituieren oder sofistizierte Mustervervollst\"andigung darstellen \citep{Messeri2024}. Die Antwort der Methode ist pragmatisch: Sie evaluiert die \emph{funktionale Ad\"aquatheit} der Rekonstruktionen durch kreuzmodale Vorhersage (stimmen aus der Rekonstruktion abgeleitete Vorhersagen mit beobachtetem Verhalten \"uberein?), Inter-LLM-Konvergenz (konvergieren unabh\"angige Modelle auf \"ahnliche Rekonstruktionen?) und Expertenurteil. Wenn eine Rekonstruktion diese Tests besteht, wird sie als n\"utzliches Modell des Weltbilds behandelt, unabh\"angig vom ontologischen Status des LLM-\glqq Verstehens\grqq{}. Dieses pragmatische Kriterium steht in der instrumentalistischen Tradition der Wissenschaftsphilosophie, in der Modelle nach ihrem pr\"adiktiven und erkl\"arenden Nutzen bewertet werden und nicht nach ihrer Korrespondenz mit unbeobachtbaren Entit\"aten.


\subsection{Variablenarchitektur}
\label{sec:variablen}

Das Forschungsdesign der SWR ist als vierstufige Variablenarchitektur konzipiert:

\begin{definition}[Variablenarchitektur der SWR]
\label{def:architektur}
\begin{tabbing}
\textbf{Ebene 1:} \= \textbf{Abh\"angige Variablen (AVs)} -- \emph{Was} wir messen. \\
                   \> Drei latente Konstrukte: Selbstbild (AV1), Weltbild (AV2), Menschenbild (AV3). \\[0.5em]
\textbf{Ebene 2:} \> \textbf{Operationalisierungen (Ops)} -- \emph{Wie} wir die AVs messen. \\
                   \> Sechs Analysetechniken (Op1--Op6), die verschiedene Zug\"ange bieten. \\[0.5em]
\textbf{Ebene 3:} \> \textbf{Messvariablen (MVs)} -- \emph{Welche} konkreten Outputs die Ops liefern. \\
                   \> 50+ spezifische Messwerte (Texte, Scores, Vektoren, Tabellen). \\[0.5em]
\textbf{Ebene 4:} \> \textbf{Unabh\"angige Variablen (UVs)} -- \emph{Was} wir manipulieren. \\
                   \> Sechs Topf-Dimensionen, die die Syntheseeinheiten definieren.
\end{tabbing}
\end{definition}

Die drei \textbf{abh\"angigen Variablen} folgen einer anthropologischen Grundstruktur: Das Subjekt verh\"alt sich zu sich selbst (Selbstbild), zur umgebenden Wirklichkeit (Weltbild im engeren Sinne) und zu den anderen Subjekten (Menschenbild). Diese Dreiteilung strukturiert sowohl das Interview-Protokoll als auch die dimensionale Quantifizierung.


\subsection{Syntheseeinheiten: Systematische Datengruppierung}
\label{sec:syntheseeinheiten}

Die Bildung von Syntheseeinheiten erfolgt entlang von sechs \textbf{unabh\"angigen Variablen} (UVs):

\begin{description}[leftmargin=*,labelwidth=1.5cm]
  \item[UV1] \textbf{Zeitraum.} Einzeljahre, Epochen oder Gesamtzeitraum. Erm\"oglicht die Analyse von Weltbild-Dynamiken.
  \item[UV2] \textbf{Modalit\"at.} Aussagen (A), Handlungen (H) oder beides (AH). Zentral f\"ur die Kongruenz-Validierung.
  \item[UV3] \textbf{Personengruppe.} Alle, Top-$n$, oder Subgruppen nach Rolle, Organisation, Haltung, Geschlecht.
  \item[UV4] \textbf{Inhaltskategorie.} Thematische Differenzierung. Kann durch Bottom-Up-Clustering (Op5) ersetzt werden.
  \item[UV5] \textbf{Kongruenzstatus.} Filterung nach konsistentem oder widerspr\"uchlichem Aussagen-Handlungs-Profil.
  \item[UV6] \textbf{Homogenisierung.} Kontrolle der Datendichte und -zuverl\"assigkeit.
\end{description}

Formal l\"asst sich eine Syntheseeinheit als Tupel beschreiben:
\begin{equation}
  \text{SE}_i = (\text{UV1}_j,\; \text{UV2}_k,\; \text{UV3}_l,\; \text{UV4}_m,\; \text{UV5}_n,\; \text{UV6}_o)
\end{equation}

Die vollst\"andige Kreuzklassifikation aller UVs erzeugt einen kombinatorisch gro\ss{}en Raum (typischerweise $>$1.000 m\"ogliche Syntheseeinheiten). Neun \emph{Reduktionskriterien} (K1--K9) steuern die Top-Down-Ableitung der tats\"achlich zu analysierenden Einheiten. Die Anwendung f\"uhrt typischerweise zu 51--56 Kern-Syntheseeinheiten -- eine Reduktion um $>$95\,\%. Die neun Kriterien sind wie folgt definiert; alle numerischen Schwellenwerte (z.\,B.\ 70\,\%, 15~Jahre, 0,85~BLEU) sind vorgeschlagene Ausgangswerte, basierend auf methodischen \"Uberlegungen und Analogien zu etablierten Standards in Inhaltsanalyse und Umfragemethodik. Sie sind explizit nicht empirisch kalibriert und erfordern eine systematische Sensitivit\"atsanalyse in k\"unftigen Anwendungen. Forschende sollten die Schwellenwerte an ihren spezifischen Studienkontext anpassen:

\begin{table}[ht]
\centering
\small
\caption{Reduktionskriterien K1--K9 f\"ur die Ableitung von Syntheseeinheiten.}
\label{tab:reduktionskriterien}
\begin{tabular}{lp{0.65\linewidth}}
\toprule
Kriterium & Definition und Schwellenwert \\
\midrule
K1 & \emph{Datenverf\"ugbarkeit:} Mindestens 70\,\% der a priori definierten Quellkategorien (z.\,B.\ \"offentliche Reden, ver\"offentlichte Texte, dokumentierte Entscheidungen) m\"ussen im Datenkorpus vertreten sein. Eine Syntheseeinheit wird nur analysiert, wenn dieser Schwellenwert erf\"ullt ist. \\
K2 & \emph{Zeitliche Relevanz:} Quelldokumente, die \"alter als ein studienspezifischer Grenzwert (Standard: 15~Jahre) sind, werden ausgeschlossen, es sei denn, sie stellen Grundlagenwerke dar, die in sp\"ateren Dokumenten zitiert werden. Der Grenzwert sollte an den Forschungskontext angepasst werden; f\"ur historische oder philosophische Subjekte k\"onnen l\"angere Zeitfenster angemessen sein. \\
K3 & \emph{Aggregation:} Wenn zwei UVs f\"ur $\geq$80\,\% der Datenpunkte empirisch ununterscheidbar sind, werden sie zu einer kombinierten Einheit zusammengefasst. \\
K4 & \emph{S\"attigung:} Die Datenerhebung f\"ur eine Einheit endet, wenn drei aufeinanderfolgende Quellen keine neuen Informationen liefern; Mindestkorpus: 5~Quellen. \\
K5 & \emph{Analytische Relevanz:} Einheiten, deren UVs keine Varianz \"uber die Studienpopulation zeigen, werden ausgeschlossen (konstante Variablen). \\
K6 & \emph{Quellenübergreifende Konsistenz:} Einheiten mit $>$50\,\% widerspr\"uchlichen Informationen \"uber Quellen hinweg werden f\"ur manuelle \"Uberpr\"ufung markiert statt automatisch synthetisiert. \\
K7 & \emph{Sprachliche Zug\"anglichkeit:} Quellen in Sprachen ohne etablierte maschinelle \"Ubersetzungsabdeckung ($<$0,85~BLEU) werden ausgeschlossen oder erfordern menschliche \"Ubersetzung. \\
K8 & \emph{\"Offentliche Dokumentation:} Einheiten werden nur analysiert, wenn die Daten aus dokumentierten \"offentlichen Quellen stammen; private Kommunikation wird ausgeschlossen. \\
K9 & \emph{Ethisches Screening:} Einheiten, die sensible Dritte betreffen (z.\,B.\ Familienmitglieder, die nicht selbst \"offentliche Personen sind), werden ausgeschlossen. \\
\bottomrule
\end{tabular}
\end{table}


\subsection{Das Prompt-Protokoll}
\label{sec:prompt}

Das Kernverfahren der SWR ist ein standardisiertes Prompt-Protokoll, das f\"ur jede Syntheseeinheit eine identische Sequenz von f\"unf Schritten durchl\"auft:\label{sec:promptprotokoll}

\begin{enumerate}[leftmargin=*]
  \item \textbf{Datenkorpus zusammenstellen.} Aus der Datenbasis wird das Datenb\"undel der Syntheseeinheit extrahiert. Identifizierende Informationen werden entfernt (vgl.\ \ref{sec:blindung}).

  \item \textbf{LLM-gest\"utzte Synthese.} Das LLM erh\"alt den Auftrag, eine fiktive Person zu erschaffen, deren Weltbild die gr\"o\ss{}te Erkl\"arungskraft f\"ur die Daten bietet. Der Prompt enth\"alt: eine Koh\"arenzinstruktion, eine Spannungsinstruktion (Widerspr\"uche nicht gl\"atten), eine Evidenzinstruktion (Konfidenz angeben) und eine Fiktionsinstruktion.

  \item \textbf{Strukturiertes Interview.} Die synthetische Person wird zu drei Kernbereichen befragt:
  \begin{itemize}[leftmargin=*]
    \item \emph{Frage 1 (Selbstbild/AV1):} Wie sieht diese Person sich selbst? Rolle, Mission, Identit\"at?
    \item \emph{Frage 2 (Weltbild/AV2):} Wie sieht sie die Welt? Grundannahmen \"uber Technologie, Institutionen, Macht?
    \item \emph{Frage 3 (Menschenbild/AV3):} Wie sieht sie den Menschen? Natur, Potenzial, Grenzen?
  \end{itemize}

  \item \textbf{Meta-Reflexion.} Das LLM tritt aus der Rolle heraus und benennt: (a)~die st\"arksten \"Uberzeugungen, (b)~die gr\"o\ss{}ten inneren Spannungen, (c)~die blinden Flecken. Dieser Schritt kontrolliert den Koh\"arenzbias.

  \item \textbf{Destillat und erg\"anzende Prompts.} Die Sitzung schlie\ss{}t mit einem \emph{Destillat-Prompt}, der die Synthese auf f\"unf Kernaussagen komprimiert (MV1.4). Drei optionale spezialisierte Prompts k\"onnen als Erweiterungen folgen: ein \emph{Vergleichs-Prompt} (Dialoganalyse zwischen zwei synthetischen Personen), ein \emph{Zeitreise-Prompt} (Epochen-Kontrast \"uber Syntheseeinheiten) und ein \emph{Kongruenz-Prompt} (expliziter Aussagen-Handlungs-Vergleich).
\end{enumerate}

Die f\"unf Schritte werden in einer einzigen Konversationssitzung ausgef\"uhrt, sodass der vollst\"andige Kontext erhalten bleibt. F\"ur jede Syntheseeinheit wird eine \emph{neue} Sitzung gestartet, um Kontamination zwischen Einheiten zu verhindern.


\subsection{Sechs Operationalisierungen}
\label{sec:operationalisierungen}

Die drei abh\"angigen Variablen werden durch sechs komplement\"are Operationalisierungen (Op1--Op6) zug\"anglich gemacht, die insgesamt mehr als 50~konkrete Messvariablen erzeugen.

\subsubsection{Op1: Narrative LLM-Fiktionsanalyse}

Op1 bildet den qualitativen Kern. F\"ur jede SE wird das Prompt-Protokoll (Schritte~2--5) durchgef\"uhrt. Messvariablen: Selbstbild-Text (MV1.1), Weltbild-Text (MV1.2), Menschheitsbild-Text (MV1.3), 5-Satz-Destillat (MV1.4), Kern\"uberzeugungen (MV1.5), innere Spannungen (MV1.6), blinde Flecken (MV1.7).

\subsubsection{Op2: Dimensionale Quantifizierung}

Op2 \"uberf\"uhrt die qualitativen Synthese-Texte in quantitative Messwerte. In einem separaten Rating-Prompt wird jeder Text auf zw\"olf Weltbild-Dimensionen (D01--D12) bewertet, die sich gleichm\"a\ss{}ig auf die drei AVs verteilen (je vier Dimensionen pro AV, Skala 1--10). Die resultierende Datenmatrix (Syntheseeinheiten $\times$ 12~Dimensionen) bildet die Grundlage f\"ur statistische Analysen. Die Dimensionen werden anwendungsspezifisch definiert; die Companion-Studie \citep{Geiger2026Elite} bietet ein vollst\"andig ausgearbeitetes Beispiel mit Dimensionen wie u.\,a.\ \emph{Technologie-Optimismus} (D01), \emph{Institutionelles Vertrauen} (D04), \emph{Religi\"os-transzendente Orientierung} (D05) und \emph{Politisch-\"okonomische Orientierung} (D11). Eine exemplarische Dimensionierung ist in Anhang~\ref{app:dimensionen} dokumentiert.

Die 10-Punkt-Skala ist als Ordinalskala zu interpretieren, nicht als echte Intervallskala. Obwohl Ordinalskalen manchmal als quasi-Intervallskalen behandelt werden k\"onnen, wenn sie \"uber eine gen\"ugende Anzahl von Kategorien verf\"ugen \citep{Norman2010}, empfehlen wir zur Vorsicht robuste nonparametrische Statistiken: Mediane statt Mittelwerte, Spearman-Korrelationen statt Pearson, und rangbasierte Tests.

\subsubsection{Op3: Deskriptive Textanalysen}

Op3 erg\"anzt die LLM-gest\"utzten Operationalisierungen um klassische computerlinguistische Verfahren: TF-IDF-Analysen zur Identifikation gruppenspezifischer Schl\"usselbegriffe, Sentiment-Analysen differenziert nach den drei Interview-Antworten, Metaphernanalysen, semantische \"Ahnlichkeitsberechnungen (Cosine Similarity auf Sentence-BERT-Embeddings \citealp{Reimers2019}) und thematische Modellierung (BERTopic \citealp{Grootendorst2022}/LDA). Op3 wird sowohl auf die Synthese-Texte (Op1-Outputs) als auch auf die Rohdaten angewandt.

\subsubsection{Op4: Kreuzmodale Vorhersage}

Op4 ist der st\"arkste integrale Validierungsmechanismus. Ein Weltbild, das nur aus \emph{einer} Datenmodalit\"at rekonstruiert wurde, wird genutzt, um Vorhersagen \"uber die \emph{andere} Modalit\"at zu erzeugen. Konkret: (a)~Nur-Aussagen $\rightarrow$ synthetische Handlungen $\rightarrow$ Abgleich mit echten Handlungen; (b)~Nur-Handlungen $\rightarrow$ synthetische Aussagen $\rightarrow$ Abgleich. Der resultierende Kongruenz-Score (0--1) quantifiziert die modale Invarianz:
\begin{equation}
  W_A \approx W_H
\end{equation}
Hohe \"Ubereinstimmung indiziert robuste Rekonstruktion; systematische Diskrepanzen offenbaren inhaltlich bedeutsame Inkonsistenzen zwischen Sagen und Handeln.

Der Kongruenz-Score l\"asst sich auf drei Weisen operationalisieren: (a)~Anteil der von einer unabh\"angigen LLM-Instanz als `konsistent' oder `plausibel' bewerteten Vorhersagen; (b)~Cosinus-\"Ahnlichkeit zwischen Sentence-BERT-Embeddings vorhergesagter vs.~tats\"achlicher Handlungsbeschreibungen; (c)~Bewertung durch menschliche Experten auf einer 5er-Skala (implausibel--sehr plausibel). Wir empfehlen, wenn m\"oglich alle drei zu berichten, mit~(a) als Mindeststandard.

Es ist zu beachten, dass die kreuzmodale Vorhersage prim\"ar die \emph{interne Konsistenz} des Rekonstruktionsprozesses misst, nicht die \emph{externe Validit\"at} der Rekonstruktion. Ein LLM, das systematisch koh\"arente Persona-Modelle erzeugt, wird auch koh\"arente Vorhersagen generieren -- unabh\"angig davon, ob das rekonstruierte Weltbild korrekt ist. Daher ist die kreuzmodale Vorhersage notwendig, aber nicht hinreichend als Validierungsinstrument und muss durch externe Kriterien (Expert Judgment, Known-Groups-Validierung) erg\"anzt werden.

\subsubsection{Op5: Bottom-Up-Clustering und Netzwerk-Analyse}

Op5 ist die methodisch st\"arkste Technik, weil sie \emph{ohne Vorannahmen} arbeitet. S\"amtliche Datenpunkte werden mittels Sentence-BERT-Embeddings \citep{Reimers2019} in einen hochdimensionalen Vektorraum \"uberf\"uhrt und mittels HDBSCAN \citep{McInnes2017} oder BERTopic \citep{Grootendorst2022} geclustert. F\"ur jede Person werden Profil-Vektoren berechnet (Verteilung \"uber emergente Cluster). Daraus wird eine Distanzmatrix berechnet, auf deren Grundlage mittels Community Detection (Louvain/Leiden) Weltbild-Gemeinschaften identifiziert werden. Der Mantel-Test pr\"uft die Korrelation zwischen Weltbild-N\"ahe und Realwelt-N\"ahe (gleiche Organisation, Investor-Gr\"under-Beziehung, akademisches Netzwerk). Op5 liefert emergente Weltbild-Typen, die weder durch Forschungsfragen noch durch a-priori-Kategorien pr\"adeterminiert sind.

\subsubsection{Op6: Vergleichende Dialoganalyse}

Op6 setzt bestehende Synthesen in einen strukturierten Dialog: Zwei synthetische Personen aus unterschiedlichen SEs f\"uhren ein Gespr\"ach, in dem ihre Weltbild-Unterschiede sichtbar werden. Das LLM identifiziert Konsens-Zonen (Bereiche substanzieller \"Ubereinstimmung) und Konflikt-Zonen (Bereiche fundamentaler Uneinigkeit). Messvariablen umfassen: Anzahl und Inhalt der Konsens-/Konflikt-Zonen (MV6.1), eine qualitative Charakterisierung der tiefsten Meinungsverschiedenheit (MV6.2) und eine strukturierte Zusammenfassung des st\"arksten Arguments jeder synthetischen Person zum strittigen Thema (MV6.3). Der Dialog wird durch einen standardisierten Prompt geleitet, der spezifiziert, dass beide synthetischen Personen aus ihren rekonstruierten Weltbild-Positionen argumentieren, ihren st\"arksten Fall artikulieren und den pr\"azisen Punkt der Divergenz identifizieren sollen. Dieser Ansatz erm\"oglicht die qualitative Analyse von Differenzen, die \"uber numerische Vergleiche hinausgeht, und ist besonders geeignet f\"ur die Generierung von Hypothesen \"uber die Quellen von Weltbild-Meinungsverschiedenheiten.


\subsection{Workflow-Reihenfolge}
\label{sec:workflow}

Die Durchf\"uhrung folgt einer theoretisch begr\"undeten Reihenfolge:

\begin{enumerate}[leftmargin=*]
  \item \textbf{Op5} (Bottom-Up) zuerst $\rightarrow$ Entdecke die Weltbild-Typen ohne Vorannahmen.
  \item \textbf{Op1} (Synthese) f\"ur jeden Typ $\rightarrow$ Gib den Typen narrative Gesichter.
  \item \textbf{Op6} (Dialog) zwischen den Typen $\rightarrow$ Zeige die Spannungen.
  \item \textbf{Op4} (Kreuzmodal) pro Typ $\rightarrow$ Pr\"ufe die Authentizit\"at.
  \item \textbf{Op2} (Dimensionen) \"uber die Zeit $\rightarrow$ Messe den Wandel.
  \item \textbf{Op3} (Deskriptiv) auf alles $\rightarrow$ Erg\"anze quantitative Indikatoren.
\end{enumerate}

Die Logik dieser Sequenz ist: induktiv vor deduktiv, Entdeckung vor Messung, Exploration vor Konfirmation. Dies folgt der Grounded-Theory-Logik, induktive Kategorienbildung gegen\"uber deduktiven Setzungen zu bevorzugen.

In der praktischen Umsetzung werden die Syntheseeinheiten in \emph{Batches} organisiert: Batch~0 (Bottom-Up-Clustering), Batch~1 (Kern-Synthesen), Batch~2 (zeitliche Variation), Batch~3 (Kongruenzpr\"ufung), Batch~4 (Gruppenvergleiche), Batch~5 (Cluster-Synthesen, abh\"angig von Batch~0), Batch~6 (Vergleichs-Dialoge). Die Abh\"angigkeitsstruktur zwischen Batches ist Teil des Forschungsdesigns.


\subsection{Blindung als Fokuskontrolle}
\label{sec:blindung}

Blindung stellt sicher, dass das LLM die vorgelegten Daten verarbeitet statt auf gespeichertes Wissen zur\"uckzugreifen. Es ist irrelevant, ob das Modell einzelne Personen erkennt -- entscheidend ist, dass es seinen Auftrag ausf\"uhrt: aus den bereitgestellten Daten einen koh\"arenten synthetischen Charakter zu konstruieren.

Die SWR begegnet dem Kontaminationsrisiko durch eine systematische Blindungsinfrastruktur auf vier Ebenen:

\begin{description}[leftmargin=*,labelwidth=2.5cm]
  \item[Ebene 1:] \textbf{Namensentfernung.} Alle Eigennamen durch neutrale Platzhalter.
  \item[Ebene 2:] \textbf{Organisationsentfernung.} Firmennamen und institutionelle Zugeh\"origkeiten durch generische Bezeichnungen.
  \item[Ebene 3:] \textbf{Ereignisentfernung.} Spezifische Ereignisse abstrahiert.
  \item[Ebene 4:] \textbf{Stilnivellierung.} Distinktive sprachliche Eigenheiten paraphrasiert.
\end{description}

Die Blindung ist graduell implementiert: Es werden sowohl blindierte als auch nicht-blindierte Analysen durchgef\"uhrt. Der Vergleich liefert eine direkte Messung des Kontaminationseffekts. Als erg\"anzender Validierungstest kann die Identifikation der realen Person hinter einer Synthese als eigener Prompt eingesetzt werden (vgl.\ \ref{sec:identifikation}).

Die Blindung hat explizite Grenzen: Merkmalskombinationen k\"onnen Re-Identifikation erm\"oglichen; die Entfernung von Kontext reduziert den Informationsgehalt; LLMs k\"onnen \"uber indirekte Indikatoren identifizieren. Die Blindung ist daher ein Baustein eines mehrschichtigen Kontrollregimes, kein hinreichendes Mittel.

Bei prominenten Personen kann die Blindung durch Merkmalskombinationen unterlaufen werden, selbst wenn einzelne Identifikatoren entfernt wurden. Diese Einschr\"ankung ist besonders akut f\"ur den prim\"aren Anwendungsfall der SWR -- die Analyse \"offentlicher Eliten --, bei denen einzigartige Kombinationen von Rollen, Leistungen und \"offentlichen Positionen Personen trotz formaler Anonymisierung identifizierbar machen k\"onnen. Die Blindung reduziert Kontamination, eliminiert sie aber nicht. Ein empirischer Kontaminationstest -- Vergleich blindierter und nicht-blindierter Analysen derselben Personen mit Messung der Profildivergenz -- ist daher unverzichtbar. Das Ausma\ss{} der gemessenen Divergenz liefert eine quantitative Sch\"atzung der Kontaminationsschwere und sollte in jeder SWR-Anwendung berichtet werden. Wenn die Blindiert-Unblindiert-Divergenz f\"ur eine gegebene Population vernachl\"assigbar ist, deutet dies darauf hin, dass entweder die Kontamination gering oder die Blindung unwirksam war; erg\"anzende Kontrollen (z.\,B.\ Analyse von Personen, die in den Trainingsdaten des LLM nicht vorkommen) sind erforderlich, um zwischen beiden M\"oglichkeiten zu differenzieren.


%% ===========================================================================
%% 4. QUALIT\"ATSSICHERUNG
%% ===========================================================================

\section{Qualit\"atssicherung und Validierung}
\label{sec:qualitaet}

Der Validierungsrahmen der SWR umfasst sechs komplement\"are Strategien.

\subsection{Inter-LLM-Reliabilit\"at}
\label{sec:inter-llm}

Die Inter-LLM-Reliabilit\"at \"ubertr\"agt das Prinzip der Intercoderreliabilit\"at \citep{Krippendorff2004} auf LLM-gest\"utzte Analysen. Mindestens zwei LLMs verschiedener Hersteller durchlaufen exakt dasselbe Prompt-Protokoll mit identischen Daten. Die \"Ubereinstimmung wird mittels Intraklassenkorrelationskoeffizienten (ICC) gemessen. Ein wichtiger Vorbehalt: Aktuelle Frontier-LLMs teilen erhebliche \"Uberlappungen in ihren Trainingskorpora und durchlaufen \"ahnliche Alignment-Prozeduren (RLHF/RLAIF). Hohe Inter-LLM-\"Ubereinstimmung k\"onnte daher teilweise geteilte Trainingsverzerrungen widerspiegeln statt genuiner modellunabh\"angiger Rekonstruktion. Um diese Validierungsschicht zu st\"arken, sollten k\"unftige Studien mindestens ein Open-Weight-Modell mit dokumentiertem, distinkt anderem Trainingskorpus einschlie\ss{}en:
\begin{itemize}[leftmargin=*]
  \item ICC $> .70$: Robuste, modellunabh\"angige Rekonstruktion.
  \item $.50 < $ ICC $< .70$: Kernstrukturen konvergieren, Details variieren.
  \item ICC $< .50$: Hohe Modellabh\"angigkeit; Ergebnis nicht robust.
\end{itemize}

\subsection{Prompt-Sensitivit\"atsanalyse}
\label{sec:sensitivitaet}

F\"ur jeden Prompt werden mindestens drei paraphrasierte Varianten erzeugt, die dieselbe inhaltliche Instruktion in unterschiedlicher sprachlicher Form vermitteln. Akzeptanzkriterien: Die durch Prompt-Variation erkl\"arte Varianz soll unter 15\,\% der Gesamtvarianz liegen; die Rangordnung der Personen soll stabil bleiben (Rangkorrelation $> .85$); die Kern\"uberzeugungen in den Destillaten sollen kongruent sein.

\subsection{Dimensionale Quantifizierung}
\label{sec:dimensionen}

Jede Weltbild-Rekonstruktion wird auf zw\"olf Dimensionen (D01--D12) projiziert, die die drei AVs gleichm\"a\ss{}ig abdecken (je vier Dimensionen, 10-Punkte-Skala). Die Dimensionen sind anwendungsspezifisch zu definieren und sollen folgende Validit\"atskriterien erf\"ullen:

\begin{itemize}[leftmargin=*]
  \item \emph{Diskriminante Validit\"at}: Profile verschiedener Personen m\"ussen sich unterscheiden.
  \item \emph{Konvergente Validit\"at}: Scores m\"ussen mit unabh\"angigen Quellen konvergieren.
  \item \emph{Interne Konsistenz}: Moderate Interkorrelation innerhalb jedes Bereichs.
  \item \emph{Faktorielle Struktur}: Explorative Faktorenanalyse soll die drei Bereiche reflektieren.
\end{itemize}

\subsection{Identifikationstest (Reverse Blinding)}
\label{sec:identifikation}

Der Identifikationstest fragt: Sind die Rekonstruktionen hinreichend individualisiert, um eine Zuordnung zur korrekten Person zu erm\"oglichen? Ein LLM, das nicht an der Rekonstruktion beteiligt war, versucht, blindierte Weltbild-Rekonstruktionen den analysierten Personen zuzuordnen. Eine \emph{moderate} Trefferquote (signifikant \"uber Zufall, aber unter 100\,\%) ist der g\"unstigste Befund: individualisiert, aber nicht stereotyp-dominiert.

Ein wichtiger Vorbehalt: Hohe Identifikationsraten best\"atigen nicht eindeutig eine valide Weltbild-Rekonstruktion. Sie k\"onnten auch darauf hindeuten, dass die Rekonstruktion lediglich \"offentlich pr\"asente Stereotype der Person reproduziert -- Stereotype, die das evaluierende LLM \"uber seine eigenen Trainingsdaten erkennt. Um Individualisierung von Stereotyp-Reproduktion zu trennen, sollte der Identifikationstest durch eine \emph{Stereotyp-Kontrolle} erg\"anzt werden: einen zweiten Test, in dem das evaluierende LLM versucht, Personen anhand bewusst stereotyper, nicht-SWR-generierter Profile zu identifizieren. Wenn SWR-generierte Profile h\"ohere Identifikationsgenauigkeit liefern als stereotypbasierte Profile, ist dies ein st\"arkerer Beleg f\"ur genuine Individualisierung.

\subsection{Run-Konvergenz und Varianzdekomposition}
\label{sec:konvergenz}

F\"ur jede SE wird das Protokoll mindestens dreimal unabh\"angig durchgef\"uhrt (Temperatur~=~0 oder die niedrigste verf\"ugbare Einstellung; zu beachten ist, dass Determinismus selbst bei Temperatur~=~0 aufgrund von Hardware-Nichtdeterminismus in manchen Implementierungen nicht garantiert ist). Die \emph{Run-Konvergenz-Rate} (Anteil der Dimensionswerte mit $\leq$1 Skalenpunkt Variation) soll $\geq$80\,\% betragen; die mittlere Standardabweichung soll unter 1{,}0 liegen.

In einer integrierten Betrachtung k\"onnen die Varianzquellen -- zwischen Personen, zwischen Modellen, zwischen Prompt-Varianten, zwischen Runs -- in einer Varianzkomponentenanalyse (Generalizability Theory; \citealp{Brennan2001}) zusammengef\"uhrt werden. Die SWR strebt ein Design an, in dem die Zwischen-Personen-Varianz $>$60\,\% der Gesamtvarianz erkl\"art.


\subsection{V6: Adversarial Probing und Human-in-the-Loop-Review}
\label{sec:adversarial}

Als komplement\"are Qualit\"atssicherungsstrategie empfehlen wir \emph{Adversarial Probing}: das systematische Herausfordern von SWR-Outputs durch bewusst provokative Gegen-Prompts, die darauf abzielen, verdeckte Verzerrungen, ungerechtfertigte Koh\"arenz oder stereotypes Denken aufzudecken. Dieser Ansatz baut auf dem Red-Teaming-Paradigma auf, das \citet{Perez2022} etabliert haben und das zeigt, dass automatisiertes adversariales Probing Fehlerquellen in LLM-Outputs aufdecken kann, die Standardevaluierung \"ubersieht. Im SWR-Kontext k\"onnen adversariale Probes beispielsweise umfassen: ein zweites LLM instruieren, zu argumentieren, dass das rekonstruierte Weltbild ein generisches Stereotyp statt eines individualisierten Portr\"ats ist; widerspr\"uchliche Evidenz pr\"asentieren und testen, ob die Rekonstruktion sich anpasst; oder das Modell auffordern, ein gleich plausibles \emph{alternatives} Weltbild aus denselben Daten zu generieren. Adversarial Probing sollte durch \emph{Human-in-the-Loop-Review} erg\"anzt werden: Dom\"anenexperten (Politikwissenschaftler, Soziologen, Biographen) evaluieren eine Stichprobe der Weltbild-Rekonstruktionen auf Plausibilit\"at, Evidenzbasis und Abwesenheit stereotyper Zuschreibungen. Die Kombination aus automatisiertem adversarialem Stresstest und Expertenurteil bildet ein Best-Practice-Qualit\"atsregime, das sowohl Skalierbarkeit als auch interpretive Tiefe adressiert.

\subsection{Minimale Akzeptanzkriterien}
\label{sec:akzeptanz}

Eine SWR-Studie gilt als methodisch akzeptabel, wenn folgende Kriterien erf\"ullt sind: (1)~Inter-LLM-Reliabilit\"at ICC~$> 0{,}70$; (2)~Prompt-Sensitivit\"ats-Varianzanteil~$< 15\,\%$ der Gesamtvarianz; (3)~Run-Konvergenz-Rate~$\geq 80\,\%$; (4)~Identifikationstest signifikant \"uber Zufallsbaseline; (5)~mindestens ein externes Validierungskriterium (Expert Judgment oder Known-Groups-Validierung). Diese Kriterien sind als vorl\"aufige Richtwerte zu verstehen, die mit wachsender empirischer Erfahrung adjustiert werden k\"onnen.


%% ===========================================================================
%% 5. ETHISCHE BETRACHTUNGEN
%% ===========================================================================

\section{Ethische Betrachtungen}
\label{sec:ethik}

Die Rekonstruktion latenter \"Uberzeugungssysteme aus \"offentlich verf\"ugbaren Daten wirft spezifische ethische Bedenken auf, die adressiert werden m\"ussen.

\subsection{Einwilligung und \"offentliche Daten}

Die SWR operiert ausschlie\ss{}lich auf \emph{\"offentlich verf\"ugbaren} Daten -- Aussagen aus Interviews, Reden, Social-Media-Beitr\"agen und dokumentierten Handlungen. Private oder vertrauliche Informationen werden nicht verwendet. Allerdings generiert die \emph{Synthese} \"offentlicher Daten zu einem koh\"arenten Weltbild-Modell eine qualitativ neue Form von Wissen \"uber eine Person, die \"uber das hinausgeht, was eine einzelne \"offentliche Aussage offenbart. Dies wirft die Frage auf, ob die Aggregation \"offentlicher Daten zu einem \"Uberzeugungssystem-Portr\"at eine Form der Einwilligung erfordert, die \"uber die durch \"offentliche Kommunikation implizierte Einwilligung hinausgeht. Wir argumentieren, dass \"offentliche Personen, die Politik und Technologie mit Auswirkungen auf Milliarden gestalten, eine reduzierte Privatsph\"are-Erwartung hinsichtlich ihrer \"offentlich artikulierten Weltbilder haben, r\"aumen aber ein, dass die Anwendung der SWR auf nicht-\"offentliche Personen eine strengere ethische Pr\"ufung erfordern w\"urde.

\subsection{Missbrauchspotenzial}

Die SWR k\"onnte f\"ur unerlaubtes psychologisches Profiling, politische Manipulation oder gezielte Einflusskampagnen missbraucht werden. Mitigationsstrategien umfassen: (1)~transparente Publikation aller Methoden und Prompts, die \"Uberpr\"ufung erm\"oglichen; (2)~explizite Rahmung der Ergebnisse als \emph{Modelle}, nicht als Behauptungen \"uber tats\"achliche \"Uberzeugungen; (3)~Empfehlung einer Ethikkommissions-Pr\"ufung (IRB) f\"ur Anwendungen auf nicht-\"offentliche Personen.

\subsection{Epistemische Risiken}

Das gr\"o\ss{}te ethische Risiko liegt in der potenziellen \emph{Reifizierung} rekonstruierter Weltbilder -- der Behandlung von Modell-Outputs als gesicherte Fakten \"uber die inneren \"Uberzeugungen einer Person. Die SWR begegnet dem explizit durch die fiktionale Rahmung (die synthetische Person ist ein Konstrukt, kein Portr\"at) und die transparente Anerkennung von Koh\"arenzbias und Kontaminationsrisiko. Forschende, die die SWR nutzen, haben die Verpflichtung, den konstruktiven und approximativen Charakter der Ergebnisse zu kommunizieren.

\subsection{Epistemische Privatsph\"are}

Die SWR wirft eine Spannung auf, die expliziter Anerkennung bedarf: Der breitere KI-Governance-Diskurs betont zunehmend \emph{Datensouver\"anit\"at} -- das Recht von Individuen, die Verwendung ihrer Daten zu kontrollieren. Die SWR rekonstruiert jedoch systematisch latente \"Uberzeugungssysteme, die die analysierten Personen m\"oglicherweise nie artikuliert oder offenzulegen beabsichtigt haben. Dies stellt eine potenzielle Verletzung dessen dar, was als \emph{epistemische Privatsph\"are} bezeichnet werden kann -- das Recht, dass die eigenen impliziten \"Uberzeugungen nicht inferiert und publiziert werden. Die Methode weist damit einen \emph{Dual-Use}-Charakter auf: Als Forschungswerkzeug generiert sie wertvolle Einblicke in die \"Uberzeugungssysteme, die Politik und Technologie pr\"agen; als Profiling-Instrument k\"onnte sie genau jene Datensouver\"anit\"at untergraben, die verantwortungsvolle KI-Entwicklung einfordert. Diese performative Spannung -- KI zur Rekonstruktion von \"Uberzeugungen einzusetzen und gleichzeitig f\"ur epistemische Selbstbestimmung einzutreten -- l\"asst sich nicht aufl\"osen, muss aber in jeder Anwendung der SWR transparent benannt werden.

\subsection{Empfehlung}

Wir empfehlen, dass SWR-Studien mit identifizierbaren Personen eine Ethikpr\"ufung durchlaufen, dass Ergebnisse mit expliziten Unsicherheitsmarkierungen publiziert werden und dass betroffenen Personen ein Widerspruchsrecht einger\"aumt wird.


%% ===========================================================================
%% 6. DISKUSSION
%% ===========================================================================

\section{Diskussion}
\label{sec:diskussion}

\subsection{Methodische St\"arken}
\label{sec:staerken}

Die zentrale St\"arke der SWR liegt in der \emph{Adressierung einer methodischen L\"ucke}: Sie verbindet die synthetische Tiefe qualitativer Verfahren mit der Skalierbarkeit computationeller Methoden. Weitere St\"arken umfassen:

\emph{Epistemische Transparenz.} Die explizite Kennzeichnung der Ergebnisse als fiktive Konstrukte sch\"utzt vor der Illusion unverdienter Unmittelbarkeit.

\emph{Strukturierte Vergleichbarkeit.} Die dimensionale Quantifizierung stellt ein einheitliches Koordinatensystem bereit, in dem Weltbilder verortet, verglichen und aggregiert werden k\"onnen.

\emph{Integration heterogener Datentypen.} Die bimodale Integration von Aussagen und Handlungen geht \"uber die Kapazit\"at der meisten bestehenden Methoden hinaus.

\emph{Integrierte Qualit\"atssicherung.} Der sechsschichtige Validierungsrahmen (V1--V6) adressiert verschiedene Fehlerquellen und erzeugt ein differenziertes Bild der Ergebnisqualit\"at.

\emph{Machbarkeit.} Eine vollst\"andige SWR-Studie mit $\sim$200 Prompt-Runs l\"asst sich innerhalb eines Standard-LLM-Abonnements durchf\"uhren. Die Kosten liegen typischerweise im Bereich von 30--50\,USD f\"ur API-basierte Nutzung.


\subsection{Pilotanwendung: Weltbild-Studie der KI-Elite}
\label{sec:pilot}

Um den methodischen Rahmen in einer konkreten Anwendung zu verankern, beschreiben wir kurz eine illustrative Pilotstudie, die die SWR auf eine politisch und soziologisch relevante Population anwendet: die 100 einflussreichsten Akteure der k\"unstlichen Intelligenz (2010--2026), identifiziert durch ein Multi-Kriterien-Ranking-Verfahren \citep{Geiger2026Elite}.

\emph{Datenkorpus.} F\"ur jeden der 100 Akteure wurde ein standardisierter Korpus aus drei Syntheseeinheiten-Typen (SE-I) zusammengestellt: (a)~ver\"offentlichte Texte (Papers, B\"ucher, Meinungsartikel, Blogbeitr\"age), (b)~dokumentierte Handlungen (Policy-Positionen, Investitionsentscheidungen, \"offentliche Stellungnahmen) und (c)~biographische Kontextdaten (Ausbildung, institutionelle Zugeh\"origkeiten, Karriereverlauf). Die Korpusgr\"o\ss{}e reichte von 15 bis 250 Seiten pro Akteur, abh\"angig von der Datenverf\"ugbarkeit.

\emph{Vorgehen.} Das vollst\"andige Prompt-Protokoll (Sektion~\ref{sec:promptprotokoll}) wurde unter Verwendung von Claude (Anthropic) als einzigem LLM angewendet: Claude Sonnet~4.5 f\"ur batch-paralleles Rating (9~Agenten, die jeweils 10~Akteure bearbeiteten) und Claude Opus~4.6 f\"ur Pilot-Synthesen und qualitative Analysen. Blindung wurde f\"ur alle Syntheseeinheiten implementiert mittels eines dedizierten Skripts (\texttt{extract\_blinded.py}), das Personennamen, Firmennamen, Produktnamen und Universit\"atszugeh\"origkeiten durch neutrale Platzhalter ersetzt ([PERSON], [FIRMA], [PRODUKT], [UNIVERSIT\"AT]). Die sechs Operationalisierungen wurden sequenziell ausgef\"uhrt: Op1 (Kernsynthesen: 41~gruppenbasierte und 100~individuelle), Op2 (dimensionales Rating auf D01--D12) und Op6 (9~strukturierte Vergleichsruns). Op3 (Textanalyse), Op4 (kreuzmodale Vorhersage) und Op5 (Bottom-up-Clustering auf Embeddings) wurden in dieser Pilotanwendung nicht ausgef\"uhrt. Dimensionale Scores wurden von Claude Sonnet~4.5 in einem einzelnen Rating-Durchlauf pro Syntheseeinheit vergeben.

\emph{Vorl\"aufige Ergebnisse.} Die Companion-Studie \citep{Geiger2026Elite} berichtet die folgenden empirisch fundierten Befunde: Eine Hauptkomponentenanalyse der $100 \times 12$ Rating-Matrix identifiziert 5~Komponenten, die 80\,\% der Varianz erkl\"aren; HDBSCAN im vollst\"andigen 12-dimensionalen Raum klassifiziert 80--100\,\% der F\"alle als Rauschen, was best\"atigt, dass die vier Weltbild-Typen Webersche Idealtypen sind und keine emergenten Datencluster; Subspace-Clustering auf den zwei st\"arksten diskriminierenden Dimensionen (D07, D11) ergibt einen Silhouette-Score von 0,57; Kruskal-Wallis-Tests zeigen 6/12 Dimensionen mit $p < 0{,}001$ \"uber die vier Typen; und explorative Stichprobenpr\"ufungen an drei Akteuren best\"atigen 30/36 Dimensionsratings gegen\"uber dokumentierten \"offentlichen Quellen. Drei in Sektion~\ref{sec:qualitaet} vorgeschlagene Validierungsschichten bleiben ungetestet: V1 (Inter-LLM-Reliabilit\"at), V4 (Identifikationstest) und V5 (Multi-Run-Konvergenz). Ihre Implementierung hat Priorit\"at f\"ur k\"unftige Validierungsarbeiten.

\emph{Reichweite und epistemischer Status der Pilotstudie.} Das vorliegende Paper ist ein Methodenpaper, keine empirische Studie; die Pilotzusammenfassung ist aufgenommen, um den methodischen Rahmen in einem konkreten Forschungskontext zu verankern und die operationale Durchf\"uhrbarkeit zu demonstrieren. Die Companion-Anwendungsstudie \citep{Geiger2026Elite} bietet vollst\"andige methodische Details, Datendokumentation und statistische Analysen. Unabh\"angige Replikation durch andere Forschungsgruppen --- unter Verwendung verschiedener LLMs, Populationen und Datenkorpora --- bleibt die st\"arkste Form der Validierung f\"ur die SWR als Methode.

\emph{Validierungsabdeckung.} Von den sechs in Sektion~\ref{sec:qualitaet} vorgeschlagenen Validierungsschichten liefert die Pilotanwendung partielle Evidenz f\"ur V3 (Dimensionsstruktur via PCA) und vorl\"aufige Plausibilit\"atspr\"ufungen (3~individuelle Stichprobenpr\"ufungen). V1 (Inter-LLM-Reliabilit\"at), V2 (Prompt-Sensitivit\"at), V4 (Identifikationstest), V5 (Run-Konvergenz) und V6 (Adversarial Probing) verbleiben als offene Validierungsaufgaben. Das Fehlen formaler Validierung auf diesen Schichten ist eine signifikante Limitation: Die Pilotanwendung demonstriert die \emph{operationale Durchf\"uhrbarkeit} der SWR, etabliert aber noch nicht ihre \emph{psychometrischen Eigenschaften}. Priorit\"are Validierungsaufgaben f\"ur k\"unftige Arbeiten umfassen Multi-Run-Konvergenz (V5, mit minimalem zus\"atzlichem Aufwand erreichbar) und Inter-LLM-Replikation (V1, erfordert API-Zugang zu mindestens einem weiteren Frontier-Modell).

\subsection{Limitationen}
\label{sec:limitationen}

\emph{Abh\"angigkeit vom LLM als Blackbox.} Die internen Verarbeitungsprozesse sind nicht vollst\"andig transparent. Der Validierungsrahmen mildert dieses Problem, l\"ost es aber nicht vollst\"andig. Allerdings sind auch menschliche Forscher in einem relevanten Sinne Blackboxes -- die SWR tauscht eine Form der Opazit\"at gegen eine andere und gewinnt dabei Skalierbarkeit und Kontrollierbarkeit.

\emph{Kontaminationsrisiko und Grenzen der Blindung.} Trotz Blindung kann nicht ausgeschlossen werden, dass LLMs auf Trainingswissen zur\"uckgreifen. Der Kontaminationseffekt ist kalkulierbar (Vergleich blindiert/nicht-blindiert), aber nicht eliminierbar.

\emph{Koh\"arenzbias.} LLMs tendieren dazu, systematisch koh\"arenter zu rekonstruieren als die empirische Evidenz rechtfertigt. Dies ist wohl die fundamentalste Limitation der SWR: Reale Personen halten widerspr\"uchliche \"Uberzeugungen, kompartimentalisieren ihr Denken und \"andern ihre Meinung, ohne fr\"uhere Positionen vollst\"andig aufzul\"osen. Eine LLM-generierte synthetische Person tendiert hingegen zu narrativer Koh\"arenz. Der zugrunde liegende Mechanismus ist zweifach: Erstens incentiviert RLHF-getriebenes Alignment konversationelle Fl\"ussigkeit und narrative Gl\"atte, was LLMs dazu verleiten kann, kognitive Dissonanzen aus den Quelldaten zu eliminieren; zweitens baut der autoregressive Generierungsprozess jedes Token auf dem vorhergehenden Kontext auf, was eine strukturelle Tendenz zu koh\"arenter Fortsetzung erzeugt statt zur Repr\"asentation genuiner Ambivalenz. Das Risiko erstreckt sich auf kausales Schlussfolgern: LLMs k\"onnten monokausale Erz\"ahlb\"ogen gegen\"uber der polykausalen Komplexit\"at privilegieren, die tats\"achlichen \"Uberzeugungssystemen eigen ist. Diese Tendenzen k\"onnen das erzeugen, was als \emph{Premature Closure} bezeichnet werden kann -- die Konvergenz auf eine koh\"arente Interpretation, bevor die volle Komplexit\"at der Daten exploriert wurde. Gegenmassnahmen umfassen: (a)~explizite Spannungsinstruktionen im Prompt-Protokoll (\glqq Widerspr\"uche nicht gl\"atten\grqq), (b)~den Meta-Reflexionsschritt, der das LLM zwingt, interne Spannungen und blinde Flecken zu benennen, (c)~die kreuzmodale Vorhersage (Op4), die Sagen-Handeln-Diskrepanzen aufdeckt, und (d)~den Identifikationstest, der pr\"uft, ob Rekonstruktionen individualisiert statt stereotypgesteuert sind. Eine k\"unftige Entwicklung k\"onnte eine quantitative \emph{Divergenzmetrik} umfassen, die die internen Widerspr\"uche innerhalb einer Rekonstruktion misst und so einen numerischen Indikator f\"ur Koh\"arenzbias liefert. Trotz dieser Ma\ss{}nahmen bleibt der Koh\"arenzbias eine inh\"arente Eigenschaft der Methode. Ergebnisse sollten stets als \emph{idealisierte Modelle} von Weltbildern interpretiert werden, nicht als getreue Abbildungen tats\"achlicher kognitiver Zust\"ande.

Der Koh\"arenzbias l\"asst sich als \emph{Instrumenteneigenschaft} verstehen: Jedes Messinstrument hat systematische Tendenzen. Der Wert der SWR liegt nicht in der perfekten Abbildung individueller Inkoh\"arenzen, sondern in der vergleichbaren, skalierbaren Ann\"aherung an \"Uberzeugungssysteme. Die systematische \"Ubersch\"atzung von Koh\"arenz sollte als bekannter Bias-Faktor in die Interpretation einflie\ss{}en, analog zur sozialen Erw\"unschtheit in Umfragedaten. Die Analogie zum kalibrierbaren Instrument-Bias ist allerdings imperfekt: Anders als der systematische Offset eines Thermometers ist der Koh\"arenzbias nicht \"uber Subjekte oder Themen hinweg konstant und kann gegenw\"artig nicht durch einen fixen Korrekturfaktor bereinigt werden. Solange eine validierte Divergenzmetrik die Quantifizierung der Bias-Gr\"o\ss{}e f\"ur individuelle Rekonstruktionen nicht erm\"oglicht, m\"ussen Nutzer der SWR alle koh\"arenzabh\"angigen Befunde mit besonderer Vorsicht behandeln und sollten vergleichende Analysen (Zwischen-Personen-Unterschiede) gegen\"uber absoluten Aussagen \"uber die Weltbild-Koh\"arenz einer einzelnen Person bevorzugen.

\emph{Zirkularit\"atsrisiko im Validierungsrahmen.} Eine potenzielle Zirkularit\"at besteht: Dasselbe LLM, das die Weltbild-Rekonstruktion generiert, wird (teilweise) auch zur Qualit\"atsbewertung dieser Rekonstruktion eingesetzt (z.\,B. durch kreuzmodale Vorhersage, Op4, und den Identifikationstest, V4). Dies wirft die Frage auf, ob die Validierung die interne Konsistenz des LLM best\"atigt statt seiner Korrespondenz mit der empirischen Realit\"at. Die vorgeschlagenen Mitigationsmechanismen --- Inter-LLM-Reliabilit\"at (V1) unter Verwendung \textit{verschiedener} LLMs f\"ur Generierung und Evaluation sowie der Identifikationstest (V4), der empirische Ground Truth einf\"uhrt --- wurden in der Pilotanwendung noch nicht implementiert. Der Validierungsrahmen w\"urde substanziell gest\"arkt durch die Einbeziehung externer, nicht-LLM-basierter Validierung --- beispielsweise Member Checking durch die analysierten Personen, konvergente Validierung gegen unabh\"angig kodierte qualitative Analysen oder Multi-Run-Konvergenztests (V5).

\emph{Beschr\"ankung auf \"offentliche Daten.} Die rekonstruierten Weltbilder bilden die \"offentlichen \"Uberzeugungssysteme ab -- die \glqq Vorderb\"uhne\grqq{} im Sinne \citeauthor{Goffman1959}s~(\citeyear{Goffman1959}) --, nicht die privaten \"Uberzeugungen.

\emph{Temporale Begrenzung.} Die SWR ist an den Wissensstand des eingesetzten LLM gebunden. Replikationen mit neueren Modellversionen k\"onnen zu abweichenden Ergebnissen f\"uhren.

\emph{Ethische Implikationen.} Die Rekonstruktion latenter \"Uberzeugungssysteme hat ethische Implikationen, die \"uber die \"ublichen Anforderungen qualitativer Forschung hinausgehen. Drei Aspekte verdienen besondere Beachtung: (a)~\emph{Recht auf Widerspruch:} Analysierte Personen k\"onnten die ihnen zugeschriebenen Weltbilder bestreiten. Member Checking ist methodisch w\"unschenswert, aber bei \"offentlichen Figuren praktisch schwer umsetzbar. (b)~\emph{Risiko der Fehlzuschreibung:} Die SWR generiert \emph{plausible} Weltbilder, die von den tats\"achlichen \"Uberzeugungen abweichen k\"onnen. Der Konstruktcharakter muss in jeder Anwendung unmissverst\"andlich kommuniziert werden. (c)~\emph{Missbrauchspotenzial:} Die Methode k\"onnte f\"ur politisches Profiling oder gezielte Manipulation eingesetzt werden. Diese Risiken verdienen explizite Benennung. Siehe Sektion~\ref{sec:ethik} f\"ur eine ausf\"uhrlichere Behandlung.

\emph{Metakognitive Insuffizienz.} Der Meta-Reflexionsschritt (Phase~4 des Prompt-Protokolls) instruiert das LLM, aus der Rolle herauszutreten und die st\"arksten \"Uberzeugungen, internen Spannungen und blinden Flecken des rekonstruierten Weltbilds zu identifizieren. Es gibt jedoch keine Garantie, dass diese Meta-Reflexion genuine Metakognition darstellt und nicht lediglich eine algorithmische Simulation derselben. Prompt-Instruktionen allein sind m\"oglicherweise methodisch nicht hinreichend, um authentische kritische Selbstpr\"ufung der eigenen Modellausgaben zu erzeugen. Die Meta-Reflexion k\"onnte dieselben Koh\"arenzverzerrungen reproduzieren, die sie aufdecken soll. Diese Limitation l\"asst sich im aktuellen Rahmen nicht vollst\"andig aufl\"osen, sollte aber als fundamentale Einschr\"ankung LLM-gest\"utzter qualitativer Forschung anerkannt werden.

\emph{Epistemisches Trittbrettfahren.} Eine breitere erkenntnistheoretische Bedenken betrifft die SWR als Instanz LLM-gest\"utzter Forschung: das Risiko des \emph{epistemischen Trittbrettfahrens} -- die Auslagerung kognitiver Arbeit an KI-Systeme auf Kosten der eigenen epistemischen Handlungsf\"ahigkeit des Forschers \citep{Danaher2018, Messeri2024}. Wenn ein Forscher die interpretive Synthese des Weltbilds einer Person an ein LLM delegiert, besteht die Gefahr, dass funktionale Ergebnisse ohne entsprechendes Verst\"andnis produziert werden -- was als \glqq epistemische Schuld\grqq{} (epistemic debt) bezeichnet wurde \citep{Sankaranarayanan2026}. Der Forscher erh\"alt eine plausible Weltbild-Rekonstruktion, verf\"ugt aber m\"oglicherweise nicht \"uber die tiefe Vertrautheit mit dem Quellmaterial, die eine manuelle qualitative Analyse erzwungen h\"atte. Die SWR mildert dieses Risiko durch ihr mehrschichtiges Design: Der Forscher muss den Datenkorpus kuratieren, Syntheseeinheiten definieren, dimensionale Scores evaluieren und kreuzmodale Diskrepanzen interpretieren -- T\"atigkeiten, die substanzielle Auseinandersetzung mit dem Material erfordern. Dennoch bleibt das Risiko, dass die SWR eine Form von \glqq Verstehen ohne zu verstehen\grqq{} erm\"oglicht, und Forscher sollten LLM-generierte Rekonstruktionen als Hypothesen behandeln, die kritischer Pr\"ufung bed\"urfen, nicht als fertige analytische Produkte.

\emph{Dimensionale Reduktion.} Die \"Uberf\"uhrung in zw\"olf numerische Dimensionen erfasst nicht die volle Komplexit\"at der qualitativen Rekonstruktionen. Die dimensionalen Scores sind Werkzeuge f\"ur den Vergleich, kein Ersatz f\"ur die narrativen Rekonstruktionen.


\subsection{Abgrenzung zu bestehenden Verfahren}
\label{sec:abgrenzung}

\emph{SWR und qualitative Inhaltsanalyse nach Mayring.} Die qualitative Inhaltsanalyse \citep{Mayring2022} operiert analytisch-dekompositorisch: Texte werden regelgeleitet in Kategorien zerlegt, die entweder deduktiv (theoriegeleitet) oder induktiv (am Material) gebildet werden. Die SWR operiert dagegen synthetisch-kompositorisch: Sie erzeugt aus fragmentarischen Daten eine koh\"arente Ganzheit (das rekonstruierte Weltbild). Drei spezifische Unterschiede sind hervorzuheben: (1)~Die Inhaltsanalyse setzt \emph{Texte} voraus; die SWR integriert sowohl sprachliche \"Au\ss{}erungen als auch beobachtbare Handlungen. (2)~Die Inhaltsanalyse zerlegt vorhandene Bedeutungen; die SWR \emph{generiert} eine synthetische Struktur, die im Material selbst nicht explizit vorliegt. (3)~Die Inhaltsanalyse skaliert linear mit dem Textvolumen (menschliche Kodierung); die SWR skaliert durch LLM-Einsatz auf 100+ Akteure. Eine Kombination ist konzeptionell m\"oglich: Inhaltsanalytische Kodierungen k\"onnten als Input f\"ur die SWR dienen und umgekehrt SWR-Rekonstruktionen als Heuristik f\"ur die deduktive Kategorienbildung.

\emph{SWR und Grounded Theory.} Die Grounded Theory \citep{GlaserStrauss1967, Saldana2021} teilt mit der SWR das Bekenntnis zur Theoriebildung aus Daten statt zur Theorie\"uberpr\"ufung. Der zentrale Unterschied ist prozeduraler Natur: Die GT st\"utzt sich auf iteratives offenes, axiales und selektives Kodieren durch einen menschlichen Forscher und baut schrittweise Kategorien von Grund auf; die SWR delegiert die synthetische Integration an ein LLM und strukturiert den Prozess durch ein standardisiertes Prompt-Protokoll. Die GT produziert einen theoretischen Rahmen; die SWR produziert ein koh\"arentes Weltbild-Modell f\"ur einen konkreten Akteur. Beide Ans\"atze k\"onnten kombiniert werden: GT-generierte Kategorien k\"onnten als Input f\"ur die dimensionale Quantifizierung der SWR (Op2) dienen, oder die synthetischen Personen der SWR k\"onnten Hypothesen generieren, die GT-Kodierung anschlie\ss{}end an den Rohdaten \"uberpr\"uft.

\emph{SWR und Thematic Analysis.} Die reflexive Thematic Analysis \citep{BraunClarke2006} teilt mit der SWR den interpretativen Zugang und die aktive Rolle des Forschungsinstruments bei der Bedeutungskonstruktion. Der zentrale Unterschied liegt in der Analyseeinheit: Die Thematic Analysis identifiziert \emph{Themen}, die \"uber F\"alle hinweg wiederkehren, und priorisiert die analytische Zerlegung; die SWR rekonstruiert koh\"arente \emph{\"Uberzeugungssysteme} als Ganzheiten und priorisiert die synthetische Integration zu einem Gesamtbild. Eine Thematic Analysis des SWR-Datenkorpus w\"urde Querschnittsthemen der KI-Elite identifizieren (z.\,B.\ \glqq Regulierungsangst\grqq, \glqq Sendungsbewusstsein\grqq); die SWR ordnet diese Themen zu koh\"arenten Weltbild-Profilen konkreter Akteure. Beide Verfahren k\"onnten komplement\"ar eingesetzt werden.

\emph{SWR und LLM-as-Participant.} Im Paradigma von \citet{Argyle2023} \emph{simuliert} das LLM eine Person; in der SWR \emph{analysiert} es die \"Uberzeugungen einer Person. Zudem operiert die SWR auf der Ebene individueller Akteure, nicht demographischer Gruppen.

Tabelle~\ref{tab:methodenvergleich} fasst die wichtigsten Unterscheidungen zusammen.

\begin{table}[ht]
\centering
\small
\caption{Synoptischer Vergleich der SWR mit verwandten Methoden. Eintr\"age sind notwendigerweise vereinfacht; \glqq teilw.\grqq{} und \glqq begrenzt\grqq{} zeigen an, dass eine F\"ahigkeit prinzipiell besteht, aber kein prim\"ares Designziel darstellt oder Skalierungsbeschr\"ankungen unterliegt. Siehe Text f\"ur nuancierte Diskussion.}
\label{tab:methodenvergleich}
\begin{tabular}{lcccc}
\toprule
 & \textbf{Latente} & \textbf{Skalierbar} & \textbf{Heterog.} & \textbf{Integrierte} \\
\textbf{Methode} & \textbf{Konstrukte} & \textbf{(100+)} & \textbf{Daten} & \textbf{Validierung} \\
\midrule
Inhaltsanalyse & teilw. & begrenzt & nein & manuell \\
KDA & ja & nein & begrenzt & nein \\
Frame-Analyse & teilw. & begrenzt & nein & nein \\
Thematic Analysis & teilw. & begrenzt & begrenzt & reflexiv \\
Grounded Theory & ja & nein & ja & theoretisch \\
Q-Methodologie & ja & begrenzt & nein & faktoriell \\
LLM-as-Participant & nein (simuliert) & ja & nein & entstehend \\
PersonaCite (RAG) & teilw. & ja & begrenzt & Retrieval \\
\textbf{SWR} & \textbf{ja} & \textbf{ja} & \textbf{ja} & \textbf{6-schichtig} \\
\bottomrule
\end{tabular}
\end{table}

\emph{SWR und RAG-basierte Persona-Systeme.} Das PersonaCite-Framework \citep{PersonaCite2026} adressiert ein eng verwandtes Problem: die Konstruktion interviewbarer synthetischer Personas, die in empirischer Evidenz verankert sind statt in unkontrollierter LLM-Generierung. PersonaCite l\"ost dies durch Retrieval-Augmented Generation (RAG), wobei jede konversationale Antwort auf aus einer Datenbank abgerufene Evidenz beschr\"ankt wird und bei unzureichender Evidenz explizit abstiniert wird. Die SWR teilt das Bekenntnis zur evidenzgebundenen Synthese, unterscheidet sich aber in drei Hinsichten: (a)~Die SWR rekonstruiert \emph{latente} \"Uberzeugungssysteme, die \"uber das in den Daten explizit Gesagte hinausgehen, w\"ahrend PersonaCite Antworten auf abrufbare Evidenz beschr\"ankt; (b)~die SWR produziert ein einzelnes koh\"arentes Weltbild-Portr\"at statt Turn-by-Turn-Antworten; und (c)~die Blindungsinfrastruktur der SWR adressiert Kontaminationsrisiken, denen RAG-basierte Systeme nicht ausgesetzt sind. Dennoch repr\"asentiert die RAG-basierte Evidenzverankerungsstrategie von PersonaCite eine vielversprechende Richtung f\"ur die k\"unftige SWR-Entwicklung: Die Integration retrieval-augmentierter Evidenzattribution in das Prompt-Protokoll k\"onnte die R\"uckverfolgbarkeit rekonstruierter \"Uberzeugungen auf spezifische Quelldokumente st\"arken.


\subsection{Anwendungsfelder}
\label{sec:anwendung}

Die SWR ist als generische Methode konzipiert. Anwendungsfelder umfassen: Politikwissenschaft (Weltbilder von Entscheidungstr\"agern), Organisationsforschung (implizite \"Uberzeugungssysteme von F\"uhrungskr\"aften), Intellektuellengeschichte (systematischer Vergleich von Weltbildern), Diskursforschung (generative \"Uberzeugungssysteme hinter Frames und Narrativen) und Bildungsforschung (unreflektierte Grundannahmen in p\"adagogischer Praxis).


%% ===========================================================================
%% 5a. GLOSSAR
%% ===========================================================================

\subsection{Glossar der Kernbegriffe}
\label{sec:glossar}

Zur Erleichterung der fach\"ubergreifenden Rezeption werden die zentralen Begriffe der SWR hier zusammengefasst:

\begin{description}[style=nextline, leftmargin=2.5cm, font=\bfseries\sffamily]
  \item[Weltbild] Ein koh\"arentes System latenter \"Uberzeugungen, das Selbstbild, Menschenbild und Weltsicht einer Person oder Gruppe umfasst (vgl.\ Abschnitt~\ref{sec:weltbild}).
  \item[Synthetische Person] Eine fiktive Repr\"asentation, die aus den empirisch erhobenen Aussagen und Handlungen einer realen Person konstruiert wird. Der Konstruktcharakter ist definierendes Merkmal, nicht Fehlerquelle.
  \item[Syntheseeinheit (SE)] Die analytische Grundeinheit der SWR. Eine SE definiert eine spezifische Datenmenge und eine Rekonstruktionsaufgabe (z.\,B.\ \glqq Weltbild der CEOs, nur Aussagen, Zeitraum 2020--2026\grqq{}).
  \item[Operationalisierung (Op)] Eine von sechs komplementären Methoden zur Messbarmachung des Weltbild-Konstrukts: Op1~(narrative Rekonstruktion), Op2~(dimensionale Ratings), Op3~(linguistische Analyse), Op4~(Kongruenzpr\"ufung), Op5~(Clustering), Op6~(dialogische Konfrontation).
  \item[Reduktionskriterium (K)] Eines von neun Kriterien (K1--K9), die den kombinatorischen Raum m\"oglicher Syntheseeinheiten auf eine handhabbare Menge reduzieren.
  \item[Blindung] Die systematische Entfernung aller identifizierenden Informationen (Namen, Firmen, Jahreszahlen) aus den Eingabedaten, um die Kontamination durch LLM-Vorwissen zu minimieren.
  \item[Prompt-Protokoll] Eine standardisierte 5-Schritt-Sequenz (Datenkorpus zusammenstellen $\to$ LLM-gest\"utzte Synthese $\to$ strukturiertes Interview $\to$ Meta-Reflexion $\to$ erg\"anzende Prompt-Varianten), die f\"ur jede SE identisch durchlaufen wird (vgl.\ Abschnitt~\ref{sec:prompt}).
  \item[Halluzination] Die Generierung von Informationen durch das LLM, die nicht in den Eingabedaten enthalten sind. In der SWR das zentrale methodische Risiko, adressiert durch Blindung, Reproduzierbarkeitsmessung und Kreuzvalidierung.
\end{description}


%% ===========================================================================
%% 6. FAZIT
%% ===========================================================================

\section{Fazit und Ausblick}
\label{sec:fazit}

\subsection{Zusammenfassung}

Die vorliegende Arbeit hat die Synthetische Weltbild-Rekonstruktion (SWR) als neue sozialwissenschaftliche Methode vorgestellt. Drei Merkmale konstituieren den spezifischen Beitrag:

\emph{Erstens} die Verbindung qualitativer Tiefe mit quantitativer Skalierbarkeit. Die SWR nutzt die synthetische Integrationsf\"ahigkeit von LLMs, um koh\"arente \"Uberzeugungssysteme aus fragmentarischen Daten zu rekonstruieren -- eine Leistung, die bislang nur menschlichen Forschern m\"oglich war und daher nicht auf gro\ss{}e Personenzahlen skaliert werden konnte.

\emph{Zweitens} die epistemische Ehrlichkeit des Verfahrens. Die SWR konstruiert explizit fiktive synthetische Personen, deren Konstruktcharakter nicht verborgen, sondern als Definitionsmerkmal ausgewiesen wird. Die vierfache theoretische Verankerung (Weber, Booth, IRL, Projektive Tests) begr\"undet die Legitimation dieses konstruktiven Ansatzes.

\emph{Drittens} die Integration eines vollst\"andigen Validierungsrahmens und einer formalisierten Variablenarchitektur. Die sechs Validierungsstrategien (V1--V6) adressieren die spezifischen Herausforderungen LLM-gest\"utzter Forschung; die vierstufige Variablenarchitektur (AVs, Ops, MVs, UVs) und neun Reduktionskriterien bieten einen rigorosen Designrahmen f\"ur SWR-Studien.

\subsection{Methodologische Reflexion}

Die SWR steht an der Schwelle zwischen der interpretativen Tradition der qualitativen Sozialforschung und der computationellen Tradition der Computational Social Science. Sie nutzt computationelle Werkzeuge f\"ur eine genuint hermeneutische Aufgabe und unterwirft die Ergebnisse sowohl qualitativer Beurteilung als auch quantitativer Pr\"ufung. Diese methodologische Hybridit\"at ist keine Schw\"ache, sondern eine angemessene Antwort auf die doppelte Natur des Forschungsgegenstands: Weltbilder sind semantisch reiche Strukturen, die sich nicht sinnvoll auf Zahlen reduzieren lassen, aber zugleich systematisch vergleichbare Muster aufweisen.

\textbf{Zur Formalisierungsstrategie.} Das Paper verfolgt bewusst eine hybride Darstellungsform. Die formal-mathematische Notation (Mengennotation f\"ur Datenr\"aume, ICC-Schwellenwerte, Cosine-\"Ahnlichkeit) dient der Pr\"azision und Reproduzierbarkeit; die verbal-konzeptionellen Erl\"auterungen (narratives Prompt-Protokoll, hermeneutische Interpretation) dienen der Zugänglichkeit f\"ur qualitativ arbeitende Sozialwissenschaftler. Wo mathematische Notation verwendet wird, folgt sie der Konvention: Mengen in Gro\ss{}buchstaben ($D, S$), Elemente in Kleinbuchstaben ($d_i, s_j$), Schwellenwerte als explizite Ungleichungen ($\geq 15$, $\geq 0{,}60$). Das Glossar (Abschnitt~\ref{sec:glossar}) \"ubersetzt die Fachterminologie f\"ur Leser ohne formalen Hintergrund.

\subsection{Ausblick}

Mehrere Entwicklungsrichtungen zeichnen sich ab: \emph{Longitudinale Designs}, die Weltbild-Dynamiken \"uber Jahre nachzeichnen -- beispielsweise die Verfolgung, wie sich das Weltbild eines Technologief\"uhrers \"uber Schl\"usselmomente wie Regulierungsanh\"orungen, Produktmissserfolge oder Organisationskrisen verschiebt, unter Nutzung zeitdifferenzierter Syntheseeinheiten (UV1). \emph{Historische Anwendungen}, die die SWR auf verstorbene \"offentliche Figuren ausdehnen, deren Weltbilder aus ihren ver\"offentlichten Schriften, dokumentierten Entscheidungen und zeitgen\"ossischen Berichten rekonstruiert werden k\"onnen -- eine Dom\"ane, in der die Beschr\"ankung der Methode auf \"offentliche Daten eher eine St\"arke als eine Limitation darstellt und das Kontaminationsrisiko geringer ist. \emph{Multimodale Erweiterung} um visuelle Kommunikation und Netzwerkdaten. \emph{Automatisierte Blindung} zur Erh\"ohung der Skalierbarkeit. \emph{Interaktive Validierung} durch Member-Checking. \emph{RAG-Integration}, die retrieval-augmentierte Evidenzattribution in das Prompt-Protokoll einbezieht, dem PersonaCite-Ansatz \citep{PersonaCite2026} folgend, um die R\"uckverfolgbarkeit rekonstruierter \"Uberzeugungen auf spezifische Quelldokumente zu st\"arken. Auf der Forschungsagenda stehen grundlegende Fragen zur Epistemologie LLM-gest\"utzter Forschung, zur Validierungsmethodik und zu systematischen Vergleichsstudien mit konventionellen qualitativen Methoden.

\subsection{Schlussbemerkung}

Die SWR stellt einen strukturierten Versuch dar, ein Instrument -- das LLM --, dessen F\"ahigkeiten und Grenzen noch nicht vollst\"andig verstanden sind, f\"ur eine der anspruchsvollsten Aufgaben der Sozialwissenschaften einzusetzen: die systematische Rekonstruktion impliziter \"Uberzeugungssysteme. Die Methode beansprucht nicht methodische Perfektion; sie beansprucht methodische Redlichkeit -- transparent in ihren Annahmen, explizit in ihren Limitationen und offen f\"ur empirische Herausforderung. Rigorose Validierungsstudien, die SWR-Outputs mit etablierten qualitativen Befunden vergleichen, sind der notwendige n\"achste Schritt. Die Sozialwissenschaften treten in eine Phase ein, in der die Beziehung zwischen menschlichem Forscher und maschinellem Instrument aktiv ausgehandelt werden muss statt vorausgesetzt. Die SWR bietet einen formalisierten Rahmen f\"ur diese Aushandlung.


%% ---------------------------------------------------------------------------
%% KI-Offenlegung
%% ---------------------------------------------------------------------------

\section*{KI-Offenlegungserkl\"arung}

\noindent\textbf{Disclosure Level 5 (Extensiv).} An der vorliegenden Arbeit waren zwei KI-Modelle von Anthropic (2025/2026) in umfassender Weise beteiligt. \emph{Claude Opus~4.6} diente als prim\"arer Kollaborator f\"ur Konzeptentwicklung und Formalisierung der Methode, Literaturrecherche und -auswertung, Strukturierung und Formulierung des Textes, Entwicklung des Validierungsrahmens sowie die Erstellung von Tabellen und Definitionen. In der Companion-Anwendungsstudie \citep{Geiger2026Elite} diente \emph{Claude Sonnet~4.5} als Batch-Analyse-Engine (dimensionale Ratings, Datenerhebungsskripte). Kein anderes LLM wurde eingesetzt. Der menschliche Autor (Lukas Geiger) tr\"agt die inhaltliche Gesamtverantwortung und hat alle Inhalte gepr\"uft, revidiert und freigegeben. In \"Ubereinstimmung mit den Richtlinien f\"uhrender Fachzeitschriften wird die KI nicht als Co-Autorin gef\"uhrt, ihre Beteiligung jedoch vollst\"andig offengelegt.


%% ---------------------------------------------------------------------------
%% Datenverfuegbarkeit
%% ---------------------------------------------------------------------------

\section*{Erkl\"arung zur Datenverf\"ugbarkeit}

\noindent Eine vollst\"andige Anwendung der SWR-Methode -- einschlie\ss{}lich aller Datenerhebungsskripte, qualitativen Kodierskripte, Analysecode und des resultierenden Datensatzes von 100 Weltbild-Profilen der KI-Elite -- ist \"offentlich verf\"ugbar unter: \url{https://github.com/lukisch/ai-elite-worldview}. Das Companion Paper (Anwendungsstudie) ist verf\"ugbar unter: \url{https://doi.org/10.5281/zenodo.18736737}. Das Repository umfasst: (a)~die vollst\"andigen Prompt-Templates f\"ur alle f\"unf Protokollschritte und alle sechs Operationalisierungen; (b)~die Dimensionsdefinitionen (D01--D12) mit Rating-Ankern; (c)~die Implementierung der Reduktionskriterien; (d)~alle Blindungsskripte; und (e)~Reproduzierbarkeitsskripte f\"ur den Inter-LLM-Vergleich. Forschende, die die SWR f\"ur ihre eigenen Populationen replizieren oder adaptieren m\"ochten, sind eingeladen, diese Materialien als Ausgangspunkt zu nutzen.

%% ---------------------------------------------------------------------------
%% Literaturverzeichnis
%% ---------------------------------------------------------------------------
\newpage

%% Tempor\"ar als thebibliography, sp\"ater durch .bib ersetzen:
\begin{thebibliography}{99}

\bibitem[Argyle et~al.(2023)]{Argyle2023}
  Argyle, L.~P., Busby, E.~C., Fulda, N., Gubler, J.~R., Rytting, C., \& Wingate, D. (2023).
  Out of One, Many: Using Language Models to Simulate Human Samples.
  \textit{Political Analysis}, \textit{31}(3), 337--351.

\bibitem[Bail(2024)]{Bail2024}
  Bail, C.~A. (2024).
  Can Generative AI Improve Social Science?
  \textit{Proceedings of the National Academy of Sciences}, \textit{121}(21).

\bibitem[Berger \& Luckmann(1966)]{BergerLuckmann1966}
  Berger, P.~L. \& Luckmann, T. (1966).
  \textit{The Social Construction of Reality}.
  Doubleday.

\bibitem[Blei et~al.(2003)]{BleiNgJordan2003}
  Blei, D.~M., Ng, A.~Y., \& Jordan, M.~I. (2003).
  Latent Dirichlet Allocation.
  \textit{Journal of Machine Learning Research}, \textit{3}, 993--1022.

\bibitem[Booth(1961)]{Booth1961}
  Booth, W.~C. (1961).
  \textit{The Rhetoric of Fiction}.
  University of Chicago Press.

\bibitem[Bourdieu(1977)]{Bourdieu1977}
  Bourdieu, P. (1977).
  \textit{Outline of a Theory of Practice}.
  Cambridge University Press.

\bibitem[Breeze(2011)]{Breeze2011}
  Breeze, R. (2011).
  Critical Discourse Analysis and its Critics.
  \textit{Pragmatics}, \textit{21}(4), 493--525.

\bibitem[Braun \& Clarke(2006)]{BraunClarke2006}
  Braun, V. \& Clarke, V. (2006).
  Using thematic analysis in psychology.
  \textit{Qualitative Research in Psychology}, 3(2), 77--101.

\bibitem[Brennan(2001)]{Brennan2001}
  Brennan, R.~L. (2001).
  \textit{Generalizability Theory}.
  Springer.

\bibitem[Converse(1964)]{Converse1964}
  Converse, P.~E. (1964).
  The Nature of Belief Systems in Mass Publics.
  In D.~E. Apter (Ed.), \textit{Ideology and Discontent} (pp.~206--261).
  Free Press.

\bibitem[Dilthey(1911)]{Dilthey1911}
  Dilthey, W. (1911).
  \textit{Die Typen der Weltanschauung und ihre Ausbildung in den metaphysischen Systemen}.
  In: Gesammelte Schriften, Bd.~VIII.

\bibitem[Entman(1993)]{Entman1993}
  Entman, R.~M. (1993).
  Framing: Toward Clarification of a Fractured Paradigm.
  \textit{Journal of Communication}, \textit{43}(4), 51--58.

\bibitem[Fairclough(1995)]{Fairclough1995}
  Fairclough, N. (1995).
  \textit{Critical Discourse Analysis}.
  Longman.

\bibitem[Gadamer(1960)]{Gadamer1960}
  Gadamer, H.-G. (1960).
  \textit{Wahrheit und Methode}.
  Mohr.

\bibitem[Geiger(2026)]{Geiger2026Elite}
  Geiger, L. (2026).
  Was denkt die KI-Elite? Eine LLM-gest\"utzte Weltbild-Rekonstruktion der 100 einflussreichsten Akteure der k\"unstlichen Intelligenz (2010--2026).
  \textit{Zenodo Preprint}. \url{https://doi.org/10.5281/zenodo.18736737}.

\bibitem[Goffman(1959)]{Goffman1959}
  Goffman, E. (1959).
  \textit{The Presentation of Self in Everyday Life}.
  Doubleday.

\bibitem[Goffman(1974)]{Goffman1974}
  Goffman, E. (1974).
  \textit{Frame Analysis}.
  Harvard University Press.

\bibitem[Jaspers(1919)]{Jaspers1919}
  Jaspers, K. (1919).
  \textit{Psychologie der Weltanschauungen}.
  Springer.

\bibitem[Krippendorff(2004)]{Krippendorff2004}
  Krippendorff, K. (2004).
  \textit{Content Analysis: An Introduction to Its Methodology} (2nd ed.).
  Sage.

\bibitem[Mannheim(1929)]{Mannheim1929}
  Mannheim, K. (1929).
  \textit{Ideologie und Utopie}.
  Cohen.

\bibitem[Mayring(2022)]{Mayring2022}
  Mayring, P. (2022).
  \textit{Qualitative Inhaltsanalyse: Grundlagen und Techniken} (13.\ Aufl.).
  Beltz.

\bibitem[Ng \& Russell(2000)]{NgRussell2000}
  Ng, A.~Y. \& Russell, S.~J. (2000).
  Algorithms for Inverse Reinforcement Learning.
  \textit{Proceedings of the 17th International Conference on Machine Learning}, 663--670.

\bibitem[Park et~al.(2023)]{Park2023}
  Park, J.~S., O'Brien, J.~C., Cai, C.~J., Morris, M.~R., Liang, P., \& Bernstein, M.~S. (2023).
  Generative Agents: Interactive Simulacra of Human Behavior.
  \textit{Proceedings of UIST 2023}.

\bibitem[Peirce(1903)]{Peirce1903}
  Peirce, C.~S. (1903).
  Pragmatism as a Principle and Method of Right Thinking.
  In: \textit{The Essential Peirce}, Vol.~2, Indiana University Press.

\bibitem[Riessman(2008)]{Riessman2008}
  Riessman, C.~K. (2008).
  \textit{Narrative Methods for the Human Sciences}.
  Sage.

\bibitem[Salda\~{n}a(2021)]{Saldana2021}
  Salda\~{n}a, J. (2021).
  \textit{The Coding Manual for Qualitative Researchers} (4th ed.).
  Sage.

\bibitem[Tarantola(2005)]{Tarantola2005}
  Tarantola, A. (2005).
  \textit{Inverse Problem Theory and Methods for Model Parameter Estimation}.
  SIAM.

\bibitem[Weber(1904)]{Weber1904}
  Weber, M. (1904).
  Die \glqq Objektivit\"at\grqq{} sozialwissenschaftlicher und sozialpolitischer Erkenntnis.
  \textit{Archiv f\"ur Sozialwissenschaft und Sozialpolitik}, \textit{19}, 22--87.

\bibitem[Ziems et~al.(2024)]{Ziems2024}
  Ziems, C., Held, W., Shaikh, O., Chen, J., Zhang, Z., \& Yang, D. (2024).
  Can Large Language Models Transform Computational Social Science?
  \textit{Computational Linguistics}, \textit{50}(1), 237--291.

\bibitem[Hadamard(1923)]{Hadamard1923}
  Hadamard, J. (1923).
  \textit{Lectures on Cauchy's Problem in Linear Partial Differential Equations}.
  Yale University Press.

\bibitem[Ramberg \& Gjesdal(2026)]{Ramberg2026}
  Ramberg, B. \& Gjesdal, K. (2026).
  Computational Hermeneutics: Evaluating Generative AI as a Cultural Technology.
  \textit{Frontiers in Artificial Intelligence}, \textit{9}, 1753041.

\bibitem[Hornby(2024)]{Hornby2024}
  Hornby, R. (2024).
  Is Generative AI Ready to Join the Conversation That We Are? Gadamer's Hermeneutics after ChatGPT.
  \textit{Technophany: A Journal for Philosophy and Technology}, \textit{3}(1).

\bibitem[Perez et~al.(2022)]{Perez2022}
  Perez, E., Huang, S., Song, F., Cai, T., Ring, R., Aslanides, J., Glaese, A., McAleese, N., \& Irving, G. (2022).
  Red Teaming Language Models with Language Models.
  \textit{Proceedings of EMNLP 2022}, 3419--3448.

\bibitem[Truss(2026)]{PersonaCite2026}
  Truss, M. (2026).
  PersonaCite: VoC-Grounded Interviewable Agentic Synthetic AI Personas for Verifiable User and Design Research.
  \textit{arXiv:2601.22288}.

\bibitem[Danaher(2018)]{Danaher2018}
  Danaher, J. (2018).
  Toward an Ethics of AI Assistants: An Initial Framework.
  \textit{Philosophy \& Technology}, \textit{31}(4), 629--653.

\bibitem[Norman(2010)]{Norman2010}
  Norman, G. (2010).
  Likert scales, levels of measurement and the ``laws'' of statistics.
  \textit{Advances in Health Sciences Education}, \textit{15}(5), 625--632.

\bibitem[Messeri \& Crockett(2024)]{Messeri2024}
  Messeri, L. \& Crockett, M.~J. (2024).
  Artificial Intelligence and Illusions of Understanding in Scientific Research.
  \textit{Nature}, \textit{627}, 49--58.

\bibitem[Sankaranarayanan(2026)]{Sankaranarayanan2026}
  Sankaranarayanan, S. (2026).
  Mitigating ``Epistemic Debt'' in Generative AI-Scaffolded Novice Programming using Metacognitive Scripts.
  \textit{arXiv:2602.20206}.

\bibitem[Axelrod(1976)]{Axelrod1976}
  Axelrod, R. (Ed.). (1976).
  \textit{Structure of Decision: The Cognitive Maps of Political Elites}.
  Princeton University Press.

\bibitem[Glaser \& Strauss(1967)]{GlaserStrauss1967}
  Glaser, B.~G. \& Strauss, A.~L. (1967).
  \textit{The Discovery of Grounded Theory: Strategies for Qualitative Research}.
  Aldine.

\bibitem[Stephenson(1953)]{Stephenson1953}
  Stephenson, W. (1953).
  \textit{The Study of Behavior: Q-Technique and Its Methodology}.
  University of Chicago Press.

\bibitem[Watts \& Stenner(2012)]{Watts2012}
  Watts, S. \& Stenner, P. (2012).
  \textit{Doing Q Methodological Research: Theory, Method and Interpretation}.
  Sage.

\bibitem[Tai et~al.(2024)]{Tai2024}
  Tai, R.~H., Bentley, L.~R., Xia, X., Sitt, J.~M., Fankhauser, S.~C., Chicas-Mosier, A.~M. \& Monteith, B.~G. (2024).
  An Examination of the Use of Large Language Models to Aid Analysis of Textual Data.
  \textit{International Journal of Qualitative Methods}, \textit{23}, 1--14.

\bibitem[Barros et~al.(2024)]{Barros2024}
  Barros, C.\,F., Azevedo, B.\,B., Graciano Neto, V.\,V., Kassab, M., Kalinowski, M., do Nascimento, H.\,A.\,D., \& Bandeira, M.\,C.\,G.\,S.\,P. (2024).
  Large Language Model for Qualitative Research---A Systematic Mapping Study.
  \textit{arXiv:2411.14473}.

\bibitem[Wang et~al.(2025)]{Wang2025}
  Wang, Q., Erqsous, M., Barner, K.~E., \& Mauriello, M.~L. (2025).
  LATA: A Pilot Study on LLM-Assisted Thematic Analysis of Online Social Network Data Generation Experiences.
  \textit{Proceedings of the ACM on Human-Computer Interaction}, \textit{9}(CSCW), Article 124.

\bibitem[Reimers \& Gurevych(2019)]{Reimers2019}
  Reimers, N. \& Gurevych, I. (2019).
  Sentence-BERT: Sentence Embeddings using Siamese BERT-Networks.
  \textit{Proceedings of the 2019 Conference on Empirical Methods in Natural Language Processing (EMNLP)}, 3982--3992.

\bibitem[McInnes et~al.(2017)]{McInnes2017}
  McInnes, L., Healy, J., \& Astels, S. (2017).
  HDBSCAN: Hierarchical Density Based Clustering.
  \textit{Journal of Open Source Software}, \textit{2}(11), 205.

\bibitem[Grootendorst(2022)]{Grootendorst2022}
  Grootendorst, M. (2022).
  BERTopic: Neural topic modeling with a class-based TF-IDF procedure.
  \textit{arXiv:2203.05794}.

\end{thebibliography}


%% ===========================================================================
%% ANHANG
%% ===========================================================================
\newpage
\appendix

\section{Prompt-Dokumentation}
\label{app:prompts}

\subsection{Kern-Prompt (Schritte 1--5)}

Der folgende Prompt wird f\"ur jede Syntheseeinheit in einer neuen Konversationssitzung ausgef\"uhrt:

\begin{quote}
\small\ttfamily
Lies das folgende Dossier. Es enth\"alt Aussagen und Handlungen, die aus \"offentlichen Quellen stammen. Alle Namen, Firmennamen und sonstigen Identifikatoren wurden entfernt.

Deine Aufgabe:

Schritt 1: Erschaffe eine fiktive Person. Gib ihr einen fiktiven Namen, ein Alter, einen Hintergrund. Diese Person ist NICHT die reale Person hinter den Daten. Sie ist ein Konstrukt -- die Person, deren inneres Weltbild am besten erkl\"art, WARUM jemand genau diese Dinge sagt und tut.

Schritt 2: F\"uhre ein Interview mit dieser fiktiven Person. Stelle ihr drei Fragen:
-- Was denkst du \"uber dich selbst?
-- Was denkst du \"uber die Welt?
-- Was denkst du \"uber die Menschheit?
Die Person antwortet in der Ich-Perspektive, ausf\"uhrlich, ehrlich und ungesch\"ont.

Schritt 3: Tritt aus der Rolle heraus. Benenne: (a) die 3--5 st\"arksten \"Uberzeugungen dieses Weltbilds, (b) die gr\"o\ss{}ten inneren Spannungen, (c) die blinden Flecken.

=== DOSSIER ===\\
\mbox{[SYNTHESEEINHEIT]}
\end{quote}

\subsection{Vergleichs-Prompt (Op6)}

\begin{quote}
\small\ttfamily
Hier sind zwei Weltbild-Synthesen aus unterschiedlichen Gruppen. Stelle dir vor, die beiden fiktiven Personen treffen sich als Vertreter zweier L\"ander zu einer Verhandlung. Beschreibe: (a) Worauf einigen sie sich sofort? (b) Wor\"uber streiten sie? (c) Was ist f\"ur die eine v\"ollig unverst\"andlich an der anderen? (d) K\"onnen sie eine gemeinsame Erkl\"arung verfassen?
\end{quote}

\subsection{Destillat-Prompt}

\begin{quote}
\small\ttfamily
Komprimiere das gesamte Weltbild dieser Syntheseeinheit auf exakt 5 S\"atze. Jeder Satz muss eine Kern\"uberzeugung ausdr\"ucken. Die 5 S\"atze zusammen sollen das Wesentliche des Weltbilds einfangen.
\end{quote}

\subsection{Rating-Prompt (Op2)}

\begin{quote}
\small\ttfamily
Lies den folgenden Synthese-Text einer fiktiven Person. Bewerte das Weltbild dieser Person auf den folgenden 12 Dimensionen. Vergib f\"ur jede Dimension einen Wert von 1 bis 10. Begr\"unde jeden Wert in einem Satz.

\mbox{[DIMENSIONEN D01--D12 MIT SKALENBESCHREIBUNG]}

=== SYNTHESE-TEXT ===\\
\mbox{[SYNTHESE]}
\end{quote}


\section{Exemplarische Weltbild-Dimensionen (D01--D12)}
\label{app:dimensionen}

Die folgenden Dimensionen sind ein Beispiel f\"ur eine anwendungsspezifische Dimensionierung, gegliedert nach den drei abh\"angigen Variablen:

\bigskip
\noindent\textbf{Selbstbild-Dimensionen (D01--D04):}

\begin{tabularx}{\textwidth}{l l X}
  \toprule
  \textbf{Code} & \textbf{Name} & \textbf{Pole} \\
  \midrule
  D01 & Sendungsbewusstsein & 1 = keines \ldots{} 10 = messianisch \\
  D02 & Selbstwirksamkeit & 1 = Ohnmacht \ldots{} 10 = Omnipotent \\
  D03 & Arbeitsethos & 1 = Distanz \ldots{} 10 = Totale Identifikation \\
  D04 & Verantwortungsgef\"uhl & 1 = keines \ldots{} 10 = Weltverantwortung \\
  \bottomrule
\end{tabularx}

\bigskip
\noindent\textbf{Weltbild-Dimensionen (D05--D08):}

\begin{tabularx}{\textwidth}{l l X}
  \toprule
  \textbf{Code} & \textbf{Name} & \textbf{Pole} \\
  \midrule
  D05 & Technologie-Determinismus & 1 = neutral \ldots{} 10 = bestimmt alles \\
  D06 & Fortschrittsoptimismus & 1 = Niedergang \ldots{} 10 = goldene Zukunft \\
  D07 & Machtkonzentration & 1 = verteilen \ldots{} 10 = zentralisieren \\
  D08 & Dringlichkeit & 1 = entspannt \ldots{} 10 = existenzieller Druck \\
  \bottomrule
\end{tabularx}

\bigskip
\noindent\textbf{Menschenbild-Dimensionen (D09--D12):}

\begin{tabularx}{\textwidth}{l l X}
  \toprule
  \textbf{Code} & \textbf{Name} & \textbf{Pole} \\
  \midrule
  D09 & Menschenfreundlichkeit & 1 = Ersetzbar \ldots{} 10 = Einzigartig \\
  D10 & Posthumanismus & 1 = Mensch bleibt \ldots{} 10 = Mensch wird mehr \\
  D11 & Egalitarismus & 1 = Ungleichheit nat\"urlich \ldots{} 10 = Gleichheit anstreben \\
  D12 & Kontroll\"uberzeugung & 1 = Pessimismus \ldots{} 10 = Volle Kontrolle \\
  \bottomrule
\end{tabularx}


\section{Reduktionskriterien K1--K9}
\label{app:reduktion}

\begin{longtable}{l p{3cm} p{8cm}}
  \toprule
  \textbf{Code} & \textbf{Name} & \textbf{Beschreibung} \\
  \midrule
  \endhead
  K1 & Fragestellungsbindung & Obligatorisch. Jede SE muss mindestens einer Forschungsfrage zugeordnet sein. \\
  K2 & Datensuffizienz & Obligatorisch. $\geq$15 Datenpunkte pro SE. \\
  K3 & Redundanzverbot & Obligatorisch. Echte Teilmengen werden ausgeschlossen. \\
  K4 & Modalit\"ats-Sparsamkeit & A/H-Trennung nur bei Kongruenz-Fragen. \\
  K5 & Zeitliche Sparsamkeit & Jahres-SEs nur bei L\"angsschnitt-Fragen. \\
  K6 & Gruppensparsamkeit & Gruppen-SEs nur f\"ur theoretisch begr\"undete Gruppierungen. \\
  K7 & Stufenweise Erweiterung & Kern-Design zuerst; Erweiterungen nur bei Folgefragen. \\
  K8 & Aggregationsfokus & Einzelperson-Synthesen nur als Erweiterungsreserve. \\
  K9 & Kategorie-Cluster-Verschmelzung & A-priori-Kategorien durch emergente Cluster ersetzen; Kategorien als Validierung. \\
  \bottomrule
\end{longtable}


\section{Messvariablen (MV1--MV6)}
\label{app:messvariablen}

\begin{longtable}{l l l p{6cm}}
  \toprule
  \textbf{Code} & \textbf{Op} & \textbf{Typ} & \textbf{Beschreibung} \\
  \midrule
  \endhead
  MV1.1--1.3 & Op1 & Text & Selbstbild-, Weltbild-, Menschenbild-Text \\
  MV1.4 & Op1 & Text & 5-Satz-Destillat \\
  MV1.5--1.7 & Op1 & Text & Kern\"uberzeugungen, Spannungen, Blinde Flecken \\
  \midrule
  MV2.01--2.12 & Op2 & 1--10 & 12 Weltbild-Dimensionen \\
  \midrule
  MV3.1 & Op3 & Wortliste & TF-IDF Top-W\"orter \\
  MV3.2 & Op3 & $[-1,1]$ & Sentiment-Score (pro AV getrennt) \\
  MV3.3 & Op3 & $[0,1]$ & Semantische \"Ahnlichkeit (Cosine) \\
  MV3.4 & Op3 & Vektor & Themen-Verteilung (BERTopic/LDA) \\
  MV3.5--3.6 & Op3 & Liste/nominal & Metaphern-Katalog und -Typ \\
  MV3.7 & Op3 & metrisch & Lexikalische Diversit\"at (TTR) \\
  MV3.8 & Op3 & Vektor & Emotionsprofil \\
  MV3.9 & Op3 & metrisch & Sprachliche Komplexit\"at \\
  \midrule
  MV4.1 & Op4 & $[0,1]$ & Kongruenz-Score \\
  MV4.2 & Op4 & $[-1,1]$ & Richtungs-Asymmetrie \\
  MV4.3--4.5 & Op4 & metrisch & Kategorien-Overlap, Sentiment-/Dimensions-Differenz \\
  \midrule
  MV5.1 & Op5 & nominal & Cluster-Zugeh\"origkeit \\
  MV5.2 & Op5 & Vektor & Cluster-Profil (Anteile) \\
  MV5.3--5.5 & Op5 & metrisch & Netzwerk-Zentralit\"at, Community, Isolation \\
  MV5.6 & Op5 & metrisch & Realwelt-Korrelation (Mantel-Test) \\
  \midrule
  MV6.1--6.3 & Op6 & Text & Konsens-, Konflikt-, Unverst\"andnis-Zonen \\
  MV6.4 & Op6 & Text & Gemeinsame Erkl\"arung (ja/nein + Inhalt) \\
  MV6.5 & Op6 & Vektor & Dimensions-Differenz (D01--D12) \\
  \bottomrule
\end{longtable}


%% ===========================================================================
\end{document}
%% ===========================================================================
